\chapter{学习范式简介}
\label{chap:machine-learning-paradigms}

第~\ref{chap:introduction}章的垃圾邮件分类器从带标签的邮件中学习：模型作出预测，再根据
标签检查哪里预测错了。如果只有很少的邮件有标签，甚至一封也没有，学习还能怎样继续？
换成下棋或控制机器人时，我们通常无法为每一种处境写出正确动作，只能观察行动的结果。
而一个长期使用的系统还会遇到新数据、新任务，以及数据无法集中到一处的限制。

这些问题都与怎样利用经验有关。学习范式涉及系统可以利用什么经验、从哪里获得反馈，以及
怎样组织学习。它与任务、模型分别回答不同的问题：任务说明要完成什么，模型说明用什么
形式表达规律，范式说明依据什么学习这些规律。
表~\ref{tab:task-paradigm-model}用同一个垃圾邮件系统区分这三个层次。

\begin{table*}[htbp]
  \centering
  \small
  \begin{tabularx}{\linewidth}{@{}lXXX@{}}
    \toprule
    \textbf{层次} & \textbf{回答的问题} & \textbf{垃圾邮件例子} & \textbf{与其他层次的关系} \\
    \midrule
    任务 & 系统要完成什么？ & 判断一封邮件是否为垃圾邮件 & 可以更换模型，任务仍是分类 \\
    模型 & 用什么形式表示输入与输出的关系？ & 逻辑回归、决策树或神经网络 & 同一模型可以用于不同学习范式 \\
    学习范式 & 系统利用什么经验和反馈？ & 从带类别标签的历史邮件学习 & 可以更换模型，仍是监督学习 \\
    \bottomrule
  \end{tabularx}
  \caption{任务、模型和学习范式描述机器学习问题的不同层次。它们相互关联，
  但不存在固定的一一对应关系。}
  \label{tab:task-paradigm-model}
\end{table*}
\FloatBarrier

同一个分类任务既可以从完整标签学习，也可以结合少量标签与大量无标签邮件，还可以先从文本
中学到有用特征，再用标签学习分类。这里沿用第~\ref{chap:introduction}章的表示概念：
用一组特征描述输入。词数、链接数等手工特征是一种表示，模型从文本
中学到的特征也可以构成表示。神经网络等模型能够利用不同来源的反馈；深度学习主要涉及模型
结构和表示学习方法，文本生成、图像生成描述任务，都不能单独确定学习范式。

本章先用监督学习、无监督学习和强化学习区分三种基本情形：已有目标可供检查、从数据中发现
结构，以及通过行动后果学习决策。有了这些基础，就能进一步讨论怎样减少对标签的依赖、提高
获取标签的效率，以及数据或任务变化后怎样共享、保留和迁移知识。学习发生的时序与协作环境
又会带来额外约束。章末把这些问题放回同一个系统，说明不同范式如何组合。

这些分组按要解决的问题组织，并非彼此互斥的分类。一种方法可以同时说明反馈从哪里来、知识
怎样复用，以及学习在哪里发生；判断一个系统采用了什么范式，需要逐项检查这些条件。

\section{监督学习}
\label{sec:supervised-learning}

如果已经积累了一批邮件及其类别，我们就有了检查预测的依据：给模型一封邮件，再用已知
类别判断它的预测是否合适。这类利用输入及其对应目标来学习预测关系的设定，称为
\term{监督学习}（supervised learning）。学习者从这些配对示例中提取规律，用于预测新输入的目标。

目标的形式取决于任务。垃圾邮件分类的目标是类别，房价回归的目标是成交价格，机器翻译的
目标则是与原文对应的译文。目标还可以是次序或其他输出序列，因此监督学习不限于分类和回归。
这些任务虽然输出不同，却都能用已有目标检查模型的预测，并据此指导学习。

“监督”并不意味着训练时始终有人站在机器旁边。手写数字的类别通常由人标注，房价预测的
目标可以来自成交记录，天气预报的目标可以来自之后实际测得的天气。这些目标来源不同，
共同点是每个训练示例都给出了当前任务希望预测的对象。目标只在训练时用于指导学习；面对
需要预测的新邮件或新房屋时，模型依靠的是已经学到的关系。

手写数字识别可以让这个过程更具体。MNIST提供手写数字图像及其类别标签，研究者用它检查
模型能否从图像与数字标签的配对中学会识别新笔迹\scite{lecun1998}。
图~\ref{fig:supervised-mnist-process}中，训练图像和标签分工不同：图像用于产生预测，标签
用于检查预测。使用阶段的新图像则直接交给训练好的模型，由模型判断它写的是哪个数字。

预测与标签的比较怎样变成学习信号？对于一个实际写着“3”的训练图像，即使模型已经把“3”排在
最可能的类别中，我们仍可以检查它给这个类别的概率。房价预测则可以比较预测价格与成交价
相差多少。损失函数把这些差异转成可计算的量，学习算法据此调整模型，使许多训练样本上的
预测更符合目标。具体调整方法可以不同，但都需要把任务要求落实为对预测的评价。

训练样本上的预测改善后，还要检查这种改善能否延续到新样本。为此，需要留出不参与训练和
调参的带标签测试样本；测试标签只用于评价，不再反过来指导模型。否则，评价的仍是模型如何
适应已经见过的答案，无法据此判断它面对新输入时的表现。

\begin{figure*}[htbp]
  \centering
  % 真矢量保留标签只到损失、反馈只回参数的监督学习设计。
\begingroup%
% 第1—3章局部矢量图元：单色黄、双色蓝红、三色蓝红黄。
\providecommand{\BasicsFiveSetup}{%
 \colorlet{B5Blue}{FigureBlue}\colorlet{B5Coral}{FigureRed}%
 \colorlet{B5Yellow}{FigureYellow}\definecolor{B5Gray}{HTML}{4C4D4F}%
 \colorlet{B5Ice}{FigureBlueFill}\colorlet{B5Rose}{FigureRedFill}%
 \colorlet{B5Cream}{FigureYellowFill}%
 \tikzset{b5 picture/.style={x=1mm,y=1mm,font=\figurefont\mdseries\fontsize{8.5}{10.5}\selectfont,text=B5Gray,line width=1pt},%
 b5 node/.style={draw=B5Gray,fill=white,line width=1pt,rounded corners=1.5mm,inner sep=3pt,font=\figurefont\mdseries\fontsize{8.5}{10.5}\selectfont,text=B5Gray,align=center},%
 b5 flow/.style={draw=B5Gray,line width=1.1pt,-{Stealth[length=2.2mm,width=1.4mm]}},%
 b5 link/.style={draw=B5Gray,line width=.7pt},%
 b5 feedback/.style={b5 flow,draw=B5Coral,dash pattern=on 2.5mm off 1.5mm}}}%
% 3—3—1网络仅为图标，不指定真实架构；各层用实心蓝、红、黄，连线中性灰。
\providecommand{\BasicsFiveNet}[3]{\begin{scope}[shift={(#1,#2)},scale=#3]%
 \foreach \a in {-3,0,3}{\foreach \b in {-3,0,3}{\draw[b5 link] (0,\a)--(5,\b);}}%
 \foreach \a in {-3,0,3}{\draw[b5 link] (5,\a)--(10,0);}%
 \foreach \a in {-3,0,3}{%
  \draw[draw=B5Gray,line width=.8pt,fill=B5Blue] (0,\a) circle (1.1);%
  \draw[draw=B5Gray,line width=.8pt,fill=B5Coral] (5,\a) circle (1.1);}%
 \draw[draw=B5Gray,line width=.8pt,fill=B5Yellow] (10,0) circle (1.1);%
\end{scope}}%
% 鸟与猫保留教学对象，不生成或模拟真实数据样本。
\providecommand{\BasicsFiveBird}[4]{\begin{scope}[shift={(#1,#2)},scale=#3]%
 \draw[draw=B5Gray,fill=#4,line width=.8pt,line join=round]%
  (-6,-1)--(-10,-2)--(-6,2).. controls (-3,3) and (-2,6) .. (1,6)%
  .. controls (4,6) and (5,4) .. (4,2).. controls (4,-1) and (-1,-3) .. (-6,-1);%
 \draw[draw=B5Gray,fill=white,line width=.7pt] (4,3)--(6,3.8)--(4,4.1)--cycle;%
 \fill[B5Gray] (2.7,4.3) circle (.35);%
 \draw[b5 link] (-5,1).. controls (-1,4) and (1,1) .. (-3,-.5);%
 \draw[b5 link] (-2,-2)--(-2,-4)--(-.7,-4);\draw[b5 link] (0,-1.5)--(0,-4)--(1.3,-4);%
\end{scope}}%
\providecommand{\BasicsFiveCat}[3]{\begin{scope}[shift={(#1,#2)},scale=#3]%
 \draw[draw=B5Gray,fill=white,line width=.8pt] (0,-1) ellipse (3.3 and 4.3);%
 \draw[draw=B5Gray,fill=white,line width=.8pt,line join=round]%
  (-3,3)--(-3,7)--(-1.2,5.8).. controls (0,6.5) and (1,6.5) .. (2,5.8)%
  --(3.3,7)--(3.3,3).. controls (2,1.3) and (-1.5,1.3) .. (-3,3);%
 \fill[B5Gray] (-1.3,4) circle (.3);\fill[B5Gray] (1.5,4) circle (.3);%
 \draw[b5 link] (0,3.4)--(0,2.7);\draw[b5 link] (-2,-2)--(-2,-5);\draw[b5 link] (2,-2)--(2,-5);%
 \draw[draw=B5Gray,line width=.8pt] (-2.5,-3).. controls (-8,-5) and (-8,-.5) .. (-5,-1);%
\end{scope}}%
\providecommand{\BasicsFiveRobot}[3]{\begin{scope}[shift={(#1,#2)},scale=#3]%
 \fill[B5Gray,rounded corners=.4mm] (-2.7,-1.6) rectangle (-1.7,1.6);%
 \fill[B5Gray,rounded corners=.4mm] (1.7,-1.6) rectangle (2.7,1.6);%
 \draw[draw=B5Gray,fill=white,line width=1pt,rounded corners=.7mm] (-1.8,-2.3) rectangle (1.8,2.3);%
 \draw[draw=B5Gray,fill=B5Yellow,line width=.8pt] (0,.6) circle (.65);%
\end{scope}}%
\BasicsFiveSetup%
\begin{tikzpicture}[b5 picture]%
 \path[use as bounding box] (0,-1) rectangle (168,79);%
 \node[anchor=west] at (0,63) {训练};\node[anchor=west] at (0,27) {使用};%
 \node[b5 node,fill=B5Ice,minimum width=25mm,minimum height=24mm] (trainx) at (16,48) {};%
 \draw[draw=B5Gray,line width=1.3pt,line cap=round] (13,53) .. controls (19,56) and (22,50) .. (16,49) .. controls (23,50) and (21,43) .. (13,45);%
 \node at (16,41) {图像 $\bm x$};%
 \node[b5 node,fill=B5Ice,minimum width=25mm,minimum height=24mm] (usex) at (16,12) {};%
 \draw[draw=B5Gray,line width=1.3pt] (16,16) ellipse (3.2 and 2.1);%
 \draw[draw=B5Gray,line width=1.3pt] (16,11.7) ellipse (3.7 and 2.5);%
 \node at (16,5) {新输入 $\bm x$};%
 \foreach \yy in {48,12}{%
  \node[b5 node,minimum width=57mm,minimum height=22mm] (model\yy) at (68,\yy) {};%
  \draw[b5 link] (67,\yy-11)--(67,\yy+11);%
  \BasicsFiveNet{45}{\yy}{1.7}%
  \node[b5 node,fill=B5Cream,minimum width=29mm,minimum height=17mm] (pred\yy) at (125,\yy) {预测 $\hat y=8$};%
  \draw[b5 flow] (model\yy.east)--(pred\yy.west);}%
 \node at (81,50) {模型 $f_{\bm\theta}$};\node at (81,44) {参数 $\bm\theta$ 待学习};%
 \node at (81,14) {模型 $f_{\bm\theta}$};\node at (81,8) {固定参数 $\bm\theta$};%
 \draw[b5 flow] (trainx.east)--(model48.west);\draw[b5 flow] (usex.east)--(model12.west);%
 \node[b5 node,fill=B5Ice,minimum width=24mm,minimum height=11mm] (label) at (155,71) {标签 $y=3$};%
 \node[draw=B5Gray,fill=B5Rose,circle,line width=1pt,minimum size=11mm,inner sep=0pt] (loss) at (155,48) {$\ell$};%
 \draw[b5 flow] (label.south)--node[right]{目标}(loss.north);%
 \draw[b5 flow] (pred48.east)--node[above]{比较}(loss.west);%
 \draw[b5 feedback] (loss.south) .. controls (154,31) and (81,28) .. (68,37);%
 \node[text=B5Coral] at (108,35) {参数更新};%
\end{tikzpicture}%
\endgroup%

  \caption{监督学习中的训练与使用。训练时图像进入模型，目标标签单独进入损失比较，虚线反馈用于更新参数；使用时只提供新图像。训练行中，数字3被暂时预测为8，原标签3只进入比较；使用行没有提供目标。图像卡与内部网络为教学示意，不指定模型架构。}
  \label{fig:supervised-mnist-process}
\end{figure*}
\FloatBarrier

监督学习的效果依赖于目标数据是否可靠、是否覆盖实际使用中的情况，以及训练数据与未来数据
是否具有足够相似的规律。若某种笔迹从未出现在训练数据中，或标签本身含有系统性错误，增加
训练轮数也不能自动解决问题。带标签数据说明了“希望预测什么”，却不保证有限样本足以代表
所有新情况。

\section{无监督学习}
\label{sec:unsupervised-learning}

现在去掉逐样本的答案。假设一批新闻没有主题标签，我们仍可以观察哪些词常一起出现、哪些
文档内容相近，借助这些规律整理新闻。这类不依赖逐样本目标答案、从数据中发现结构或分布
规律的设定，称为
\term{无监督学习}（unsupervised learning）。

要寻找哪一种规律，取决于数据的用途。\term{聚类}（clustering）把相似样本组织成群组，
例如根据购买记录探索顾客的行为模式；\term{降维}（dimensionality reduction）用更少的
变量保留数据中的重要变化或邻近关系，便于可视化、压缩或后续分析；\term{分布建模}
（distribution modeling）描述哪些观测或组合较常出现，可以用于生成样本，或为识别不常见的
观测提供依据。没有标签时，这些目标仍然可以用于检查学习结果，只是检查的依据来自数据本身。

以聚类为例，一种思路是反复调整分组，使组内样本在选定的表示下更相似。若只使用顾客的购买
金额，分组可能主要反映消费水平；加入购买时间和商品种类，结果又可能反映购物习惯。
同一批顾客可以有不同分组，因为我们选择了不同的表示和相似性标准。“无监督”省去了逐样本的
答案，并没有省去学习目标与假设。

回到新闻整理的问题，一篇报道可能既涉及体育，也涉及经济，强行把它放入一个群组会丢掉部分
信息。潜在狄利克雷分配（latent Dirichlet allocation, LDA）采用另一种描述：一篇文档可以
混合多个主题，每个主题用词语出现的概率来刻画。原始研究曾在
16,000篇TREC AP新闻文档上学习100个主题，并观察各主题中概率较高的词\scite{blei2003}。
训练数据没有为每篇新闻预先填写主题标签；学习过程根据词语的共现规律，也就是哪些词经常
出现在同一篇文档中，形成主题。研究者随后根据代表性词语解释其含义。
图~\ref{fig:unsupervised-topic-process}展示了从文档集合到主题表示的
过程：训练时利用整批文档中的统计规律，使用时则可以据此描述一篇新文档的主题组成。

\begin{figure}[!htb]
  \centering
  \bookdraftgraphic{figures/scenes/v1r5-unsupervised.png}
  \caption{无监督主题发现的过程示意。没有主题标签的新闻集合提供词语共现信息，学习过程形成
  若干潜在主题；同一篇文档可以由多个主题共同构成，主题名称是在结果解释阶段由人给出的。}
  \label{fig:unsupervised-topic-process}
\end{figure}
\FloatBarrier

主题发现与监督式新闻分类因此有所不同。监督式分类预先规定“体育”“财经”等类别，并用
示例标签引导学习；主题发现先寻找词语共同出现的模式，再由人解释结果。即使人最后为某个
主题取名为“体育”，这个名字也不是训练时提供的监督标签。一篇报道同时涉及经济和体育，
还可以在主题表示中保留多种成分。

从数据中发现的结构，不一定直接对应人关心的概念。例如，把顾客的几十项购买特征压缩成
两个数，就可以把每位顾客画成平面上的一个点，观察哪些点聚在一起。但这两个数是降维方法
计算出的坐标，横轴和纵轴未必分别对应“收入”“年龄”等容易解释的属性；图上分开的群组，
也不一定就是业务上需要区分的顾客类型。

因此，发现结构之后，还需要结合具体用途检查它是否有意义。以新闻主题为例，可以查看一个
主题中的代表词和文档是否表达了相近内容，换一批相似新闻重新学习时主要主题是否仍然存在，
以及用这些主题整理文档后是否更便于检索。这些检查分别帮助判断结果能否解释、是否稳定，
以及是否有实际用途，不能仅凭一张分组清晰的图就认定学习结果有效。

\section{强化学习}
\label{sec:reinforcement-learning}

监督学习和无监督学习的例子都从已有数据出发。下棋时，情况更进一步：系统当前选择哪一步，会
改变棋盘，也会改变之后需要面对的选择。机器人移动、车辆换道同样如此。我们不仅要判断
当前行动是否合适，还要考虑它给后续行动留下了什么处境。

\term{强化学习}（reinforcement learning）利用行动带来的奖励反馈学习决策，目标通常是
提高预期累计回报，即综合多步奖励，并考虑各种可能交互结果的平均表现。作出决策的学习者称为
\term{智能体}（agent），与它交互的外部世界称为\term{环境}（environment）。例如，
下棋程序是智能体，棋局及对手的回应构成它面对的环境；\term{奖励}（reward）则是环境
提供的数值评价，用来说明某次交互的结果有多好。

智能体需要根据当前可得的信息选择行动。这条选择规则称为\term{策略}（policy）。环境的
\term{状态}（state）描述当前处境，\term{观测}（observation）则是智能体实际拿到的信息。
两者未必相同：机器人摄像头只能拍到某个视角，未必能看到所有障碍物。因此，决策有时还要
结合先前的观测和行动记录，不能把眼前一帧画面当作完整状态。

奖励与监督标签不同。标签可以直接告诉分类器某张图片是什么，奖励通常只评价行动造成的结果，
不会说明当时应当选择哪个动作；奖励还可能延迟出现。例如，下围棋时胜负要等一局结束才能
确定，眼前看似有利的落子却可能导致后面输棋。若只按眼前吃子多少来评价，就可能鼓励不利于
最终获胜的选择。因此，强化学习需要把当前行动与后续结果联系起来。

记策略为$\pi$，长度为$T$的交互轨迹为$\tau$，第$t$步奖励为$r_t$。策略与环境共同
诱导轨迹分布$p_\pi$。一种常见目标是最大化期望折扣回报
\begin{equation}
  j_{\gamma}(\pi)
  =\mathbb{E}_{\tau\sim p_\pi}
    \left[\sum_{t=0}^{T-1}\gamma^t r_t\right],
  \qquad 0\leq\gamma\leq1.
  \label{eq:paradigms-expected-return}
\end{equation}
求和先评价一条轨迹，外层期望再按各种可能轨迹出现的概率取平均。参数$\gamma$称为
\term{折扣因子}（discount factor）。有限任务取$\gamma=1$时得到未折扣总回报；
取$\gamma<1$会降低远期奖励的权重，在奖励有界的无限持续过程中也能使折扣和保持有限。
折扣因子因而既影响计算，也改变不同时间奖励的相对权重。

对于长期持续运行的系统，也可以在极限存在时评价单位时间的平均回报：
\begin{equation}
  j_{\mathrm{avg}}(\pi)
  =\lim_{T\to\infty}\frac1T
    \mathbb{E}_{\tau\sim p_\pi}
    \left[\sum_{t=0}^{T-1}r_t\right].
  \label{eq:paradigms-average-return}
\end{equation}
折扣回报与平均回报对应不同评价口径；仅比较期望还可能掩盖结果波动，有限交互对长期表现的
估计也可能不稳定\scite{sutton2018}。

在边交互边学习的设定中，智能体按当前策略行动，记录观测、动作、奖励与后续观测。学习
方法利用这些记录评价长期效果，再调整策略，让可能带来较好回报的行动更容易被选择。
策略一变，下一轮遇到的情形也会改变。如果只重复已知有效的行动，可能错过更好的选择；
如果不断尝试未知行动，又可能放弃已经能获得的收益。这就是探索与利用之间的权衡
\scite{sutton2018}。

怎样根据经验改进策略，有几种常见思路。\emph{基于价值的方法}估计处于某个状态或采取某个
行动后能获得的长期回报，再据此选择行动。\emph{直接优化策略的方法}则直接调整行动规则，
以提高预期回报。\term{行动者—评论家}（actor--critic）方法结合了两者：行动者选择动作，
评论家评估这些选择，为行动规则的更新提供依据。

如果还能够预测环境怎样变化，就可以在真正行动前比较不同选择的后果。为此可以使用已知
或从数据中学到的\term{环境模型}（environment model），即描述行动如何影响后续状态或
观测及奖励的模型，再借助它进行规划。无论采用哪种思路，都需要把眼前反馈与长期结果联系起来。

2015年的深度Q网络（deep Q-network, DQN）研究让智能体仅根据Atari游戏画面和游戏分数
学习控制，在49款游戏上使用相同的基本算法与模型设定\scite{mnih2015}。
智能体看到画面后选择按键，游戏环境更新画面和分数；这些交互不断产生新的经验，智能体据此
改善后续行动。图~\ref{fig:reinforcement-atari-process}用机器人探索场景说明同样的交互分工：
行动改变处境，新观测与奖励提供后续学习的依据。深度Q网络属于基于价值的方法：它学习估计
不同按键在后续游戏中可能带来的
累计回报，而不只是预测按下之后立刻增加多少分。这里的游戏分数提供奖励信息，但不会直接
给出“此刻应按哪个键”。

\begin{figure*}[t]
  \centering
  % 真矢量导航场景与交互架构，压紧高度以便正文回填。
\begingroup%
% 第1—3章局部矢量图元：单色黄、双色蓝红、三色蓝红黄。
\providecommand{\BasicsFiveSetup}{%
 \colorlet{B5Blue}{FigureBlue}\colorlet{B5Coral}{FigureRed}%
 \colorlet{B5Yellow}{FigureYellow}\definecolor{B5Gray}{HTML}{4C4D4F}%
 \colorlet{B5Ice}{FigureBlueFill}\colorlet{B5Rose}{FigureRedFill}%
 \colorlet{B5Cream}{FigureYellowFill}%
 \tikzset{b5 picture/.style={x=1mm,y=1mm,font=\figurefont\mdseries\fontsize{8.5}{10.5}\selectfont,text=B5Gray,line width=1pt},%
 b5 node/.style={draw=B5Gray,fill=white,line width=1pt,rounded corners=1.5mm,inner sep=3pt,font=\figurefont\mdseries\fontsize{8.5}{10.5}\selectfont,text=B5Gray,align=center},%
 b5 flow/.style={draw=B5Gray,line width=1.1pt,-{Stealth[length=2.2mm,width=1.4mm]}},%
 b5 link/.style={draw=B5Gray,line width=.7pt},%
 b5 feedback/.style={b5 flow,draw=B5Coral,dash pattern=on 2.5mm off 1.5mm}}}%
% 3—3—1网络仅为图标，不指定真实架构；各层用实心蓝、红、黄，连线中性灰。
\providecommand{\BasicsFiveNet}[3]{\begin{scope}[shift={(#1,#2)},scale=#3]%
 \foreach \a in {-3,0,3}{\foreach \b in {-3,0,3}{\draw[b5 link] (0,\a)--(5,\b);}}%
 \foreach \a in {-3,0,3}{\draw[b5 link] (5,\a)--(10,0);}%
 \foreach \a in {-3,0,3}{%
  \draw[draw=B5Gray,line width=.8pt,fill=B5Blue] (0,\a) circle (1.1);%
  \draw[draw=B5Gray,line width=.8pt,fill=B5Coral] (5,\a) circle (1.1);}%
 \draw[draw=B5Gray,line width=.8pt,fill=B5Yellow] (10,0) circle (1.1);%
\end{scope}}%
% 鸟与猫保留教学对象，不生成或模拟真实数据样本。
\providecommand{\BasicsFiveBird}[4]{\begin{scope}[shift={(#1,#2)},scale=#3]%
 \draw[draw=B5Gray,fill=#4,line width=.8pt,line join=round]%
  (-6,-1)--(-10,-2)--(-6,2).. controls (-3,3) and (-2,6) .. (1,6)%
  .. controls (4,6) and (5,4) .. (4,2).. controls (4,-1) and (-1,-3) .. (-6,-1);%
 \draw[draw=B5Gray,fill=white,line width=.7pt] (4,3)--(6,3.8)--(4,4.1)--cycle;%
 \fill[B5Gray] (2.7,4.3) circle (.35);%
 \draw[b5 link] (-5,1).. controls (-1,4) and (1,1) .. (-3,-.5);%
 \draw[b5 link] (-2,-2)--(-2,-4)--(-.7,-4);\draw[b5 link] (0,-1.5)--(0,-4)--(1.3,-4);%
\end{scope}}%
\providecommand{\BasicsFiveCat}[3]{\begin{scope}[shift={(#1,#2)},scale=#3]%
 \draw[draw=B5Gray,fill=white,line width=.8pt] (0,-1) ellipse (3.3 and 4.3);%
 \draw[draw=B5Gray,fill=white,line width=.8pt,line join=round]%
  (-3,3)--(-3,7)--(-1.2,5.8).. controls (0,6.5) and (1,6.5) .. (2,5.8)%
  --(3.3,7)--(3.3,3).. controls (2,1.3) and (-1.5,1.3) .. (-3,3);%
 \fill[B5Gray] (-1.3,4) circle (.3);\fill[B5Gray] (1.5,4) circle (.3);%
 \draw[b5 link] (0,3.4)--(0,2.7);\draw[b5 link] (-2,-2)--(-2,-5);\draw[b5 link] (2,-2)--(2,-5);%
 \draw[draw=B5Gray,line width=.8pt] (-2.5,-3).. controls (-8,-5) and (-8,-.5) .. (-5,-1);%
\end{scope}}%
\providecommand{\BasicsFiveRobot}[3]{\begin{scope}[shift={(#1,#2)},scale=#3]%
 \fill[B5Gray,rounded corners=.4mm] (-2.7,-1.6) rectangle (-1.7,1.6);%
 \fill[B5Gray,rounded corners=.4mm] (1.7,-1.6) rectangle (2.7,1.6);%
 \draw[draw=B5Gray,fill=white,line width=1pt,rounded corners=.7mm] (-1.8,-2.3) rectangle (1.8,2.3);%
 \draw[draw=B5Gray,fill=B5Yellow,line width=.8pt] (0,.6) circle (.65);%
\end{scope}}%
\BasicsFiveSetup%
\begin{tikzpicture}[b5 picture]%
 \path[use as bounding box] (0,-13) rectangle (168,69);%
 \foreach \left in {0,63,126}{%
  \begin{scope}[shift={(\left,34)}]%
   \draw[draw=B5Gray,fill=white,line width=1pt] (0,0) rectangle (42,34);%
   \foreach \x/\y/\w/\h in {2/26/20/6,2/12/5/12,15/17/8/6,32/12/5/13,36/5/4/4,16/2/12/4}{%
    \draw[draw=B5Gray,fill=B5Ice,line width=.8pt] (\x,\y) rectangle (\x+\w,\y+\h);}%
   \draw[draw=B5Coral,fill=B5Rose,line width=1pt] (37,30) circle (2.2);%
   \draw[draw=B5Blue,line width=.8pt,dashed] (7,8) .. controls (14,12) and (18,12) .. (23,14) .. controls (25,15) and (25,18) .. (26,20) .. controls (27,26) and (33,28) .. (37,30);%
  \end{scope}}%
 \fill[B5Ice] (7,42)--(13.9,47.8)--(8.6,50.9)--cycle;\BasicsFiveRobot{7}{42}{1}%
 \fill[B5Ice] (86,48)--(94.5,50.1)--(89.8,56.2)--cycle;\BasicsFiveRobot{86}{48}{1}%
 \draw[draw=B5Coral,dashed,line width=.8pt] (163,64) circle (4.3);\BasicsFiveRobot{163}{64}{1}%
 \draw[b5 flow] (45,51) .. controls (51,54) and (55,54) .. (60,51);%
 \draw[b5 flow] (108,51) .. controls (114,54) and (118,54) .. (123,51);%
 \node at (21,29) {局部观测 $o_t$};\node at (84,29) {移动后 $o_{t+1}$};\node at (147,29) {到达目标};%
 \node[b5 node,minimum width=27mm,minimum height=13mm,fill=B5Cream] (policy) at (18,15) {策略\\$h_t$};%
 \node[b5 node,minimum width=37mm,minimum height=13mm,fill=B5Ice] (world) at (83,15) {$s_t\longrightarrow s_{t+1}$};%
 \node[b5 node,minimum width=38mm,minimum height=16mm,fill=B5Rose] (learn) at (143,15) {学习\\$(o_t,a_t,r_t,o_{t+1})$};%
 \draw[b5 flow] (policy.east)--node[above]{$a_t$}(world.west);%
 \draw[b5 flow] (world.east)--node[above]{记录}(learn.west);%
 \draw[b5 flow,rounded corners=2mm] (world.south)--(83,1)--node[below]{$o_{t+1}$}(12,1)--([xshift=-6mm]policy.south);%
 \draw[b5 feedback,rounded corners=2mm] (learn.south)--(143,-9)--node[below]{更新策略}(24,-9)--([xshift=6mm]policy.south);%
\end{tikzpicture}%
\endgroup%

  \caption{机器人探索的概念场景与强化学习的交互分工。上部导航示意展示局部观测、移动与到达目标，不表示测量结果或精确状态转移。下部路径表示策略、环境和学习之间的分工：策略依据历史选择动作，学习利用含奖励的交互记录更新策略；非终止时继续交互。奖励不是正确动作标签，到达目标后的回合终止只属于上部教学场景。}
  \label{fig:reinforcement-atari-process}
\end{figure*}

在前面的围棋例子中，棋盘局面是状态，落子是行动，胜负在一局结束后提供关键反馈。程序还
可以通过自我对弈产生新的棋局，让不同局面的后续胜负为学习提供经验。
第~\ref{sec:combining-learning-paradigms}节将以AlphaGo说明这些经验怎样与人类棋谱结合。

在机器人控制中，观测可以来自摄像头、关节位置和力传感器，行动可以是移动关节、调整抓取姿态
或选择行走方向，奖励则用于表达是否抓稳物体、是否到达目标以及动作是否平稳。与可以反复
重开的游戏不同，真实机器人一次失败就可能损坏设备，经验的获取本身也有成本。

自动驾驶中的路径规划、跟车和换道也可以写成交互决策问题：车辆根据周围交通选择动作，之后的
位置、速度和其他车辆反应形成新状态，安全、通行效率与乘坐舒适度共同影响评价。不过，自动
驾驶系统包含多个学习问题：环境感知可以采用监督或自监督学习，决策模块则需要考虑行动
的后果。类似的序贯决策还出现在推荐、资源调度和能源管理中，是否适合使用强化学习取决于
问题需要怎样描述行动的后果、反馈和长期目标。

有些决策问题可以暂时不考虑行动对后续环境状态的影响。例如，推荐一条新闻后，系统只记录
这次是否点击，并不建模它对用户后续兴趣的影响。但系统仍然只知道被展示新闻的结果，不知道
同一时刻改推其他新闻会怎样。这类只观察所选行动反馈、通过反复选择积累奖励的问题，
称为\term{多臂老虎机问题}（multi-armed bandit）。若选择前还能观察用户特征、候选新闻等
信息，就得到\term{上下文老虎机问题}（contextual bandit）。Li等人的新闻推荐研究采用了
后一设定\scite{li2010bandit}。

老虎机问题是强化学习中一种简化的决策设定：当前行动不通过改变环境状态影响后续奖励，
但它带来的信息仍会影响以后的选择，所以仍需平衡探索与利用。若一次推荐还会改变用户的后续
行为，并且这些变化是优化目标的一部分，就需要更完整的序贯决策模型。区别在于对行动后果的
建模范围，不能仅凭系统用于推荐就确定设定。

这些例子带出一个实际问题：如果不能承受反复试错的代价，交互经验从哪里来？可以先在模拟器
中学习，或用专家示范提供较可靠的初始行为；也可以采用\term{离线强化学习}
（offline reinforcement learning），在学习阶段
利用已经收集的交互记录，而不再额外与环境交互\scite{levine2020offline}。这些办法可以减少
在真实环境中重新试错的次数，但模拟器与真实环境之间仍可能存在差异，历史记录也未必覆盖将来会遇到的
状态和动作。因此，学到的策略仍需要经过检验，再逐步用于实际环境。

即使经验足够，奖励也未必完整表达真正的目标。如果只奖励机器人尽快搬完物体，它可能学会快速搬运，
却忽略物体是否受损。针对这一问题，可以把任务完成情况、损坏程度等共同纳入评价，对不能
接受的行为设置额外约束；对于难以直接写成数值的要求，还可以请人比较不同的行为结果，
从偏好反馈中学习评价信号，第~\ref{sec:combining-learning-paradigms}节会介绍这种组合。
无论采用哪种方式，都需要检查实际行为是否符合任务要求，不能只以训练奖励是否提高判断成功。

专家示范还提供了另一种反馈：直接展示在某种处境下应该怎样行动。从这些示范中学习行动方式，
称为\term{模仿学习}（imitation learning）。其中的\term{行为克隆}（behavior cloning）
把专家遇到的观测作为输入、专家动作作为目标，形成一个监督学习问题。

行为克隆的难处在于，学习者的一次失误可能让它进入示范没有覆盖的状态。例如，一步走偏后，
后面的观测就不再与原来的路线相同，照搬已有示范可能继续出错。DAgger等方法因而让学习者
在环境中行动，再请专家为它实际遇到的状态提供动作示范，补上原先没有覆盖的情形
\scite{ross2011}。模仿学习和强化学习都能处理连续决策，但前者主要从示范学习行动，后者
主要根据行动结果的奖励改进策略。仅凭系统会连续行动，还不能确定它采用了哪种范式。

除了直接学习专家动作，还可以尝试从示范中推断行为背后的目标。\term{逆强化学习}
（inverse reinforcement learning）尝试寻找能够解释专家行为的奖励函数，再利用推断出的
奖励学习策略\scite{ng2000irl}。例如，示范路线不仅提供“在此处左转”的动作，还可能帮助
推断专家对路程、障碍和安全距离的取舍。不过，同一组行为可能符合多种奖励，有限示范通常
不能唯一确定专家的真实目标；从示范学习奖励后，也仍需检验由它引导的行为。

\section{知识的高效利用}
\label{sec:efficient-knowledge-use-overview}

在监督学习中，标签告诉模型需要区分什么，但获取可靠标签往往比收集原始数据更费力。
例如，照片可以批量采集，确认每张照片中的物种却可能需要专家。若标注预算有限，就要考虑
怎样让已有标签发挥更大作用、利用其余无标签数据，或把新增标注用在更有价值的样本上。

本节围绕这几种办法讨论知识的高效利用。“高效”主要指减少对任务标签的依赖，或提高标签的
获取和使用效率，并不意味着计算量一定更小。无标签数据中的结构、已有规则和样本间的关系
都可能提供额外依据，但能否帮助目标任务，仍需用可靠的评价来判断。

\subsection{半监督学习}
\label{sec:semi-supervised-learning}

商品识别中的类别可能已经确定，却来不及给全部图片标注。少量带标签
图片可以说明类别含义，其余图片则可能补充不同光照、
角度和背景下的样貌。\term{半监督学习}（semi-supervised learning）把这两部分数据一起
用于目标任务，让标签提供方向，让无标签数据补充对输入规律的认识。

不过，没有答案的图片怎样指导模型？一种办法是利用模型当前已经学到的判断。
\term{自训练}（self-training）先用有标签数据训练模型，再让它预测无标签样本。
其中被采纳为临时训练目标的预测，称为\term{伪标签}（pseudo-label）。模型用原有标签与
选中的伪标签继续学习，之后还可以重新生成伪标签。这里把预测当作临时目标，是为了利用更多
数据，并不意味着这些预测已经变成正确答案。

筛选伪标签时，常会参考模型对预测的确定程度，即\term{置信度}（confidence）。但模型可能
很确定地犯错；如果错误伪标签被反复用于训练，原来的误判反而会得到强化。因此，除了筛选，
还需要控制伪标签对学习的影响。

另一种办法无需先认定图片属于哪一类，而是检查模型对同一对象的判断是否一致。例如，轻微
改变亮度通常不应改变图片中动物的类别。要求这些变换前后的预测保持稳定，称为
\term{一致性约束}（consistency regularization）。同一原图经不同变换得到的版本称为
不同\term{视图}（view）；即使没有类别标签，模型也能比较各视图的预测，把分歧作为调整依据。
这个要求依赖变换保留任务含义。若裁剪已经去掉了待识别的动物，再要求类别不变就可能误导学习。

相似样本之间的联系也可以提供线索。\term{基于图的方法}（graph-based methods）把样本
画成节点，把相似样本连起来，再利用少量已知标签向邻近的无标签节点传播类别信息。经典的
图半监督学习用两类样本共同构图，约束相邻节点具有相近预测\scite{zhu2003}。
相关的聚类思路则假设同一密集群组往往属于同一类别。无论怎样组织相似关系，都需要检查它
是否与任务有关：如果图片只是因为背景颜色相近而被连在一起，传播的类别就未必可靠。

这些办法也可以共同使用。FixMatch将伪标签与一致性约束结合：模型对轻微变换后的无标签
图片作出预测，若置信度达到要求，就把预测当作临时目标，要求同一图片经过更强变换后仍给出
相同判断\scite{sohn2020}。真实标签负责约束类别含义，伪标签和一致性要求则把学习扩展到
没有提供标签的图片上。

在包含十类物体图片的CIFAR-10实验中，研究者只向模型提供少量训练图片的类别标签，其余
训练图片按上述方式参与学习。这里的“无标签”是实验时隐藏了答案，模型不能在训练中查看
这些真实类别。图~\ref{fig:semi-supervised-image-process}展示了两路数据怎样共同形成分类能力。

\begin{figure}[!htb]
  \centering
  \bookdraftgraphic{figures/scenes/v1r5-semi-supervised.png}
  \caption{半监督学习的过程示意。少量带类别标签的图片规定目标，大量无标签图片提供额外结构
  信息，两路经验共同用于形成分类能力；模型图标内的实心点与空心环分别示意两路学习信号，汇入同一分类器，不表示精确网络结构。无标签图片的隐藏真值不能作为训练答案使用。}
  \label{fig:semi-supervised-image-process}
\end{figure}
\FloatBarrier

在这个例子中，无标签图片仍来自目标任务需要识别的十个类别。如果额外加入大量不属于这些
类别的图片，模型也可能自信地把它们判断为某个已知类别，再把错误判断用作伪标签。因此，
使用额外图片时，需要检查它们是否与目标任务相关，不能只看数量是否增加。

利用无标签数据时，还要明确最终需要预测哪些样本。如果目标是学到一条能够处理以后新样本的
规则，就属于\term{归纳学习}（inductive learning）。若待预测的一批图片已经给定，学习
可以利用这些图片本身的结构，目标是为这批特定图片补全答案，则属于\term{直推学习}
（transductive learning），也称转导学习。直推支持向量机就是利用给定待分类样本的例子
\scite{joachims1999transductive}。

两者区别在于预测范围与学习时可见的输入。半监督描述的是有标签与无标签数据共同参与学习，
既可以采用归纳设定，也可以采用直推设定。直推学习允许利用预先给定的待预测输入，但这些
样本的真实标签仍需隐藏；它在这批样本上的表现，也不能直接当成对未来新样本的评价。

半监督学习让已有类别标签与无标签图片共同服务于目标任务。若暂时没有这些类别标签，
还可以考虑从原始数据中构造另一个有答案的学习问题，先积累可复用的知识。

\subsection{自监督学习}
\label{sec:self-supervised-learning}

标注稀缺时，还有一条思路：原始数据能不能同时提供问题和答案？例如，一段完整文本没有
情感类别标签，但如果暂时遮住其中一个词，原文就保留着恢复这个词所需的答案。我们无需另请
标注者逐句填写，就能构造大量“根据上下文补全内容”的训练示例。

\term{自监督学习}（self-supervised learning）就是从数据自身的内容或关系中构造训练目标。
除遮挡恢复外，还可以根据前文预测后续内容，或学习匹配同一对象的不同视图。“自”指目标
来自数据，并不表示模型自行决定学什么；遮住哪些内容、比较哪些对象，都属于目标设计。
设计这些问题，是希望模型为完成预测而学到有用规律：补词可能需要语法与语义，恢复图像区域
可能需要纹理、形状及物体关系。模型究竟学到了多少，还要检查它在实际任务中的表现。

BERT采用的掩码预测展示了怎样构造这类示例\scite{devlin2019}。处理文本时，系统先将句子分成
\term{词元}（token），即模型逐个处理的文本单位，可以是一个字、一个词或词的一部分。
训练时随机选取一部分词元，把其中多数替换为表示隐藏内容的特殊标记，即掩码；模型再根据
左右上下文预测原词元，原始文本中被遮住的内容就是答案。
同一句话因此同时提供输入和目标。例如，在教学句子“猫在垫子上”中
暂时隐藏“垫子”对应的词元，输入是保留其余内容的句子，目标是被隐藏的原词元。
图~\ref{fig:self-supervised-masked-process}展示了这个例子中的两路信息：可见内容用于预测，
保留的原词用于检查预测是否正确。

\begin{figure}[htbp]
  \centering
  % 例字8.5pt，文楷与框线同次矢量排版；答案不进入模型。
\begingroup%
% 第1—3章局部矢量图元：单色黄、双色蓝红、三色蓝红黄。
\providecommand{\BasicsFiveSetup}{%
 \colorlet{B5Blue}{FigureBlue}\colorlet{B5Coral}{FigureRed}%
 \colorlet{B5Yellow}{FigureYellow}\definecolor{B5Gray}{HTML}{4C4D4F}%
 \colorlet{B5Ice}{FigureBlueFill}\colorlet{B5Rose}{FigureRedFill}%
 \colorlet{B5Cream}{FigureYellowFill}%
 \tikzset{b5 picture/.style={x=1mm,y=1mm,font=\figurefont\mdseries\fontsize{8.5}{10.5}\selectfont,text=B5Gray,line width=1pt},%
 b5 node/.style={draw=B5Gray,fill=white,line width=1pt,rounded corners=1.5mm,inner sep=3pt,font=\figurefont\mdseries\fontsize{8.5}{10.5}\selectfont,text=B5Gray,align=center},%
 b5 flow/.style={draw=B5Gray,line width=1.1pt,-{Stealth[length=2.2mm,width=1.4mm]}},%
 b5 link/.style={draw=B5Gray,line width=.7pt},%
 b5 feedback/.style={b5 flow,draw=B5Coral,dash pattern=on 2.5mm off 1.5mm}}}%
% 3—3—1网络仅为图标，不指定真实架构；各层用实心蓝、红、黄，连线中性灰。
\providecommand{\BasicsFiveNet}[3]{\begin{scope}[shift={(#1,#2)},scale=#3]%
 \foreach \a in {-3,0,3}{\foreach \b in {-3,0,3}{\draw[b5 link] (0,\a)--(5,\b);}}%
 \foreach \a in {-3,0,3}{\draw[b5 link] (5,\a)--(10,0);}%
 \foreach \a in {-3,0,3}{%
  \draw[draw=B5Gray,line width=.8pt,fill=B5Blue] (0,\a) circle (1.1);%
  \draw[draw=B5Gray,line width=.8pt,fill=B5Coral] (5,\a) circle (1.1);}%
 \draw[draw=B5Gray,line width=.8pt,fill=B5Yellow] (10,0) circle (1.1);%
\end{scope}}%
% 鸟与猫保留教学对象，不生成或模拟真实数据样本。
\providecommand{\BasicsFiveBird}[4]{\begin{scope}[shift={(#1,#2)},scale=#3]%
 \draw[draw=B5Gray,fill=#4,line width=.8pt,line join=round]%
  (-6,-1)--(-10,-2)--(-6,2).. controls (-3,3) and (-2,6) .. (1,6)%
  .. controls (4,6) and (5,4) .. (4,2).. controls (4,-1) and (-1,-3) .. (-6,-1);%
 \draw[draw=B5Gray,fill=white,line width=.7pt] (4,3)--(6,3.8)--(4,4.1)--cycle;%
 \fill[B5Gray] (2.7,4.3) circle (.35);%
 \draw[b5 link] (-5,1).. controls (-1,4) and (1,1) .. (-3,-.5);%
 \draw[b5 link] (-2,-2)--(-2,-4)--(-.7,-4);\draw[b5 link] (0,-1.5)--(0,-4)--(1.3,-4);%
\end{scope}}%
\providecommand{\BasicsFiveCat}[3]{\begin{scope}[shift={(#1,#2)},scale=#3]%
 \draw[draw=B5Gray,fill=white,line width=.8pt] (0,-1) ellipse (3.3 and 4.3);%
 \draw[draw=B5Gray,fill=white,line width=.8pt,line join=round]%
  (-3,3)--(-3,7)--(-1.2,5.8).. controls (0,6.5) and (1,6.5) .. (2,5.8)%
  --(3.3,7)--(3.3,3).. controls (2,1.3) and (-1.5,1.3) .. (-3,3);%
 \fill[B5Gray] (-1.3,4) circle (.3);\fill[B5Gray] (1.5,4) circle (.3);%
 \draw[b5 link] (0,3.4)--(0,2.7);\draw[b5 link] (-2,-2)--(-2,-5);\draw[b5 link] (2,-2)--(2,-5);%
 \draw[draw=B5Gray,line width=.8pt] (-2.5,-3).. controls (-8,-5) and (-8,-.5) .. (-5,-1);%
\end{scope}}%
\providecommand{\BasicsFiveRobot}[3]{\begin{scope}[shift={(#1,#2)},scale=#3]%
 \fill[B5Gray,rounded corners=.4mm] (-2.7,-1.6) rectangle (-1.7,1.6);%
 \fill[B5Gray,rounded corners=.4mm] (1.7,-1.6) rectangle (2.7,1.6);%
 \draw[draw=B5Gray,fill=white,line width=1pt,rounded corners=.7mm] (-1.8,-2.3) rectangle (1.8,2.3);%
 \draw[draw=B5Gray,fill=B5Yellow,line width=.8pt] (0,.6) circle (.65);%
\end{scope}}%
\BasicsFiveSetup%
\begin{tikzpicture}[b5 picture]%
 \path[use as bounding box] (0,-10) rectangle (111,71);%
 \foreach \xx/\word in {7/猫,20/在,46/上}{%
  \node[b5 node,minimum width=12mm,minimum height=9mm] at (\xx,51) {\word};%
  \node[b5 node,minimum width=12mm,minimum height=9mm] at (\xx,32) {\word};}%
 \node[b5 node,minimum width=14mm,minimum height=9mm,fill=B5Ice] (original) at (33,51) {垫子};%
 \node[b5 node,minimum width=14mm,minimum height=9mm,fill=B5Cream] (masked) at (33,32) {[MASK]};%
 \node[b5 node,minimum width=24mm,minimum height=10mm,fill=B5Rose] (target) at (99,51) {垫子：$y$};%
 \draw[b5 flow] (original.north) .. controls (44,68) and (86,68) .. (target.north);%
 \node at (65,68) {复制};%
 \draw[b5 flow] (original.south)--node[right]{遮蔽}(masked.north);%
 \draw[b5 link,rounded corners=1.5mm] (1,26)--(1,24)--(52,24)--(52,26);%
 \node[b5 node,minimum width=40mm,minimum height=18mm] (model) at (29,10) {};%
 \BasicsFiveNet{14.5}{10}{1.35}\node at (39.5,10) {$f_{\bm\theta}$};%
 \draw[b5 flow] (29,24)--(model.north);%
 \node[b5 node,minimum width=13mm,minimum height=10mm] (pred) at (76,10) {$\hat y$};%
 \node[draw=B5Gray,fill=B5Rose,circle,line width=1pt,minimum size=11mm,inner sep=0pt] (loss) at (99,10) {$\ell$};%
 \draw[b5 flow] (model.east)--(pred.west);\draw[b5 flow] (pred.east)--(loss.west);%
 \draw[b5 flow] (target.south)--(loss.north);%
 \draw[b5 feedback,rounded corners=2mm] (loss.south)--(99,-6)--(29,-6)--(model.south);%
 \node[text=B5Coral] at (67,-3) {参数更新};%
\end{tikzpicture}%
\endgroup%

  \caption{掩码预测以“猫在垫子上”为教学例子。输入将“垫子”替换为[MASK]，原词单独保留为比较目标；模型只能读可见上下文，不能读被保留的答案。括线汇合全部可见词元，珊瑚红虚线将损失反馈到参数；内部网络只作模型图元。}
  \label{fig:self-supervised-masked-process}
\end{figure}
\FloatBarrier

这样得到输入与目标后，就能像监督学习一样比较预测和答案，再调整模型。改用下一词元预测时，
输入只保留目标之前的内容，以原文中紧接其后的词元为目标，同样能自动构造大量训练配对。
因此，理解自监督学习
的关键是追溯目标怎样产生：掩码预测的答案是暂时藏起的原文，监督式情感分类的答案则是
另行提供的正面或负面类别。差别在于训练问题的构造方式，不能只用标签是否由人手工填写来判断。

自监督学习与无监督学习有共同动机：利用没有人工任务标签的数据。仍以同一批新闻为例，前面的
主题发现直接从词语共现中学习潜在主题；掩码预测则先隐藏部分词元，建立“根据上下文恢复原词”
的预测问题。前者主要用学习到的结构描述数据，后者显式地把数据转成可比较预测与目标的训练
示例。两者都有学习目标，也都包含人为设定，不能把差别理解为“有没有目标函数”。

在广义用法中，自监督学习可以纳入无监督学习，因为它同样不依赖额外的人工任务标签；在强调
训练信号的来源时，人们又常把它单独讨论。本书采用后一种组织方式，以便说明数据怎样被转化
为训练目标。

自监督目标可以直接服务于文本生成，也可以先帮助模型学习通用表示，再把这些表示用于分类、
问答等具体任务。这种先积累可复用知识的阶段称为\term{预训练}（pre-training），后续要
完成的具体任务称为\term{下游任务}（downstream task）。例如，先从大量文本中学习词语与
上下文的关系，再用少量带情感标签的评论学习区分正面和负面评价。

预训练不会自动给出下游任务的类别含义，最终任务仍可能需要标签或其他说明。自监督学习能否
帮助下游任务，取决于构造的目标是否促使模型学到有用规律。例如，若模型只靠相邻词的搭配就能
补全被遮住的内容，它未必因此学会理解整段文字。

\subsection{弱监督学习}
\label{sec:weakly-supervised-learning}

从数据自身构造目标是一条途径，利用成本较低但质量较弱的现有标注是另一条途径。
要学习定位图像中的汽车，理想的标注应当告诉模型汽车在哪里。但手头也许只有“这张图含有
汽车”的说明；从网页收集的图片甚至可能带着错误标签。\term{弱监督学习}
（weakly supervised learning）研究怎样利用这类不够完整、精确或可靠的信号学习，而不是
要求所有数据都先达到理想标注标准。

一种有代表性的划分据此区分三种困难：只有部分样本有标签，叫作\emph{监督不完整}
（incomplete supervision）；只有整图类别而缺少任务所需的位置等信息，叫作
\emph{监督不精确}（inexact supervision）；标签与实际内容不符，叫作
\emph{监督不准确}（inaccurate supervision）\scite{zhou2018}。第~\ref{sec:semi-supervised-learning}节
的半监督学习主要处理监督不完整，因而也属于这一广义弱监督范围。本节重点看粗标签、
错误标签，以及如何从规则中获得监督。

如果标注者能说出一些判断经验，还可以把经验写成规则，让规则为许多样本产生弱标签。
Snorkel支持这种做法：领域专家编写\term{标注函数}（labeling function），也就是根据
关键词、模式或外部知识给样本一个判断的程序。它可以在依据不足时“弃权”，不作判断。
例如，垃圾评论识别可以把“含有链接”作为一条可疑线索，但正常分享也会附带链接，单条规则
不一定可靠。Snorkel综合多条可能一致、冲突或弃权的规则，估计样本属于各类别的概率，
再用这些概率标签训练最终模型\scite{ratner2017}。图~\ref{fig:weak-supervision-process}
中的汇合环节，正是把各条规则的判断转成可供模型使用的监督。

用规则提供弱标签可以减少逐条标注的工作，也会把规则的局限带入学习过程。例如，多条规则若都
把含链接的评论判为垃圾内容，综合之后仍可能误判正常的分享链接。粗粒度标签也可能遗漏任务
需要的信息。因此，评价最终模型时还需要独立的可靠标签；只检查它是否符合原有规则，无法
判断它是否学会了真正需要的分类能力。

\begin{figure}[htbp]
  \centering
  % 同列为同一评论，同行为同一规则；箭头文字位于路径上方。
\begingroup%
\begin{tikzpicture}[foundation picture]
 \def\BPositive{正}\def\BNegative{负}
 \path[use as bounding box] (0,-25) rectangle (111,35);
 \foreach \x/\name in {22/评论1,36/评论2,50/评论3}{\node at (\x,29) {\name};}
 \foreach \y/\name/\a/\b/\c in {15/规则1/正/正/弃,0/规则2/正/负/负,-15/规则3/弃/负/正}{
  \node[anchor=east] at (14,\y) {\name};
  \foreach \x/\v in {22/\a,36/\b,50/\c}{
   \edef\BVote{\v}\def\BVoteFill{white}%
   \ifx\BVote\BPositive\def\BVoteFill{FigureBlueFill}\fi%
   \ifx\BVote\BNegative\def\BVoteFill{FigureRedFill}\fi%
   \node[foundation model,fill=\BVoteFill,circle,minimum size=8mm,inner sep=0pt] at (\x,\y) {\v};
  }
 }
 \draw[FigureGrid,line width=.8pt] (16,22)--(56,22);
 \node[foundation model,minimum width=15mm,minimum height=19mm,align=center] (merge) at (75,0) {规则\\汇合};
 \draw[foundation flow] (56,0)--node[above,yshift=2mm]{综合}(merge.west);
 \draw[foundation flow] (merge.east)--node[above,yshift=2mm]{监督}(93,0);
 \draw[FigureStroke,line width=1.1pt] (93,-15)--(93,15);
 \foreach \y in {-15,0,15}{\draw[foundation flow] (93,\y)--(97.5,\y);}
 \foreach \y/\i in {15/1,0/2,-15/3}{
  \node[foundation node,fill=FigureYellowFill,minimum width=13mm,minimum height=9mm] at (104,\y) {$\tilde y_\i$};
 }
\end{tikzpicture}
\endgroup%

  \caption{三条规则在三个未标注评论上的教学投票。正、负、弃分别表示两种倾向与弃票；同一列可能存在冲突。综合规则信号形成训练监督$\tilde y_i$，图中不规定合并算法，也不把某条规则输出等同于真实标签；评价仍使用独立可靠标签。}
  \label{fig:weak-supervision-process}
\end{figure}
\FloatBarrier

回到缺少位置标注的汽车图像，粗标签与细预测之间需要有一条合理的联系。可以把图像看作多个
候选区域的集合：既然整图含有汽车，就至少有一个区域应当包含汽车，而不要求每个区域都是
汽车。这是\term{多示例学习}（multiple-instance learning）的一种典型解释：监督约束的
是一组对象，而不是逐一给出组内
每个对象的答案。把整图标签直接赋给所有区域，会把背景也错误地当作汽车。

面对\emph{带噪声标签}（noisy labels），常见思路是识别可疑样本、降低其影响，或显式描述
标签可能怎样出错，再据此调整学习。例如，可以借助少量经过核验的样本估计标签的可靠性。
这里不能把所有难学的样本都视为错标：少见但正确的情形也可能让模型犯错。处理弱监督的
核心，是把已知的不确定性纳入学习过程，而不是无条件地把现有标签当作真值。

对于Snorkel这类多来源规则监督，通常先估计弱标签来源的可靠性并综合它们的判断，再让最终模型
从原始样本及综合后的监督信号中学习。综合过程不能简单地把每一条规则都看成独立的一票；
两条仅关键词略有不同的规则可能提供了几乎相同的证据。合理的方法需要考虑来源的准确程度、
覆盖范围以及可能的依赖关系，并允许保留不确定性。

这也解释了为什么产生弱标签之后还要训练模型：规则先覆盖一部分典型情形，模型再从原始内容
中学习可推广的预测规律，尝试处理规则没有覆盖的新样本。规则提供了起点，但其中的系统性
错误也可能被模型继承。弱监督的效果最终要由独立样本上的可靠评价来判断。

\subsection{对比学习}
\label{sec:contrastive-learning}

除了恢复被遮住的内容，数据中已知的对应关系也能提供学习依据。
一张照片轻微改变颜色或取景后，我们仍希望模型认出同一个对象。即使没有类别标签，两张图
来自同一原图这一关系也已知，可以用来指导表示学习。\term{对比学习}（contrastive learning）
通过匹配与不匹配的关系学习表示，常见做法是让匹配样本的表示接近，同时让不匹配样本相对分离。
希望匹配的一对样本称为\term{正样本对}（positive pair），用于对照的一对不匹配样本称为
\term{负样本对}（negative pair）。这里的“正、负”描述样本之间的关系，与二分类中的
正类和负类含义不同。

SimCLR希望先从大量没有提供类别标签的图片中学习视觉表示，再用这些表示完成图像分类。
原研究在ImageNet等数据上先进行不使用类别标签的表示学习，之后再用
带标签数据训练分类器或调整模型，检验这些表示能否帮助识别图像中的对象\scite{chen2020}。

SimCLR把表示学习转化成“找出同一张原图的另一个视图”的问题。每次取一批原始图片，对每张图片
分别进行两组随机变换，例如裁剪和颜色变化，得到两个视图；再让同一个模型把各视图转换成
可比较的数值表示。对于其中一个视图，来自同一原图的另一个视图是匹配对象，其余原图产生的
视图则提供对照。只要求匹配视图接近还不够，因为把所有图片都变成同一个表示也能满足这一点。
加入其他图片作对照，才要求模型既保留同一对象在变换下的联系，又保留不同图片之间的区别。

图~\ref{fig:contrastive-learning-process}用同一原图的两个变换视图说明这个过程：取景和颜色
可以不同，但它们来自同一原图；另一幅猫图的视图作为对照。模型在这一阶段并不知道
“鸟”“猫”这些类别，反馈只来自视图是否由同一原图生成，因此这里属于自监督学习。
反复比较不同图片后，模型有机会学到不随这些变换轻易改变的视觉特征，再将其用于后续分类。
图中的几何位置只表示匹配视图应当接近，不是实验中计算得到的表示分布。

\begin{figure*}[htbp]
  \centering
  % 同源鸟视图与他源猫的三个输入共享同一编码器参数。
\begingroup%
% 第1—3章局部矢量图元：单色黄、双色蓝红、三色蓝红黄。
\providecommand{\BasicsFiveSetup}{%
 \colorlet{B5Blue}{FigureBlue}\colorlet{B5Coral}{FigureRed}%
 \colorlet{B5Yellow}{FigureYellow}\definecolor{B5Gray}{HTML}{4C4D4F}%
 \colorlet{B5Ice}{FigureBlueFill}\colorlet{B5Rose}{FigureRedFill}%
 \colorlet{B5Cream}{FigureYellowFill}%
 \tikzset{b5 picture/.style={x=1mm,y=1mm,font=\figurefont\mdseries\fontsize{8.5}{10.5}\selectfont,text=B5Gray,line width=1pt},%
 b5 node/.style={draw=B5Gray,fill=white,line width=1pt,rounded corners=1.5mm,inner sep=3pt,font=\figurefont\mdseries\fontsize{8.5}{10.5}\selectfont,text=B5Gray,align=center},%
 b5 flow/.style={draw=B5Gray,line width=1.1pt,-{Stealth[length=2.2mm,width=1.4mm]}},%
 b5 link/.style={draw=B5Gray,line width=.7pt},%
 b5 feedback/.style={b5 flow,draw=B5Coral,dash pattern=on 2.5mm off 1.5mm}}}%
% 3—3—1网络仅为图标，不指定真实架构；各层用实心蓝、红、黄，连线中性灰。
\providecommand{\BasicsFiveNet}[3]{\begin{scope}[shift={(#1,#2)},scale=#3]%
 \foreach \a in {-3,0,3}{\foreach \b in {-3,0,3}{\draw[b5 link] (0,\a)--(5,\b);}}%
 \foreach \a in {-3,0,3}{\draw[b5 link] (5,\a)--(10,0);}%
 \foreach \a in {-3,0,3}{%
  \draw[draw=B5Gray,line width=.8pt,fill=B5Blue] (0,\a) circle (1.1);%
  \draw[draw=B5Gray,line width=.8pt,fill=B5Coral] (5,\a) circle (1.1);}%
 \draw[draw=B5Gray,line width=.8pt,fill=B5Yellow] (10,0) circle (1.1);%
\end{scope}}%
% 鸟与猫保留教学对象，不生成或模拟真实数据样本。
\providecommand{\BasicsFiveBird}[4]{\begin{scope}[shift={(#1,#2)},scale=#3]%
 \draw[draw=B5Gray,fill=#4,line width=.8pt,line join=round]%
  (-6,-1)--(-10,-2)--(-6,2).. controls (-3,3) and (-2,6) .. (1,6)%
  .. controls (4,6) and (5,4) .. (4,2).. controls (4,-1) and (-1,-3) .. (-6,-1);%
 \draw[draw=B5Gray,fill=white,line width=.7pt] (4,3)--(6,3.8)--(4,4.1)--cycle;%
 \fill[B5Gray] (2.7,4.3) circle (.35);%
 \draw[b5 link] (-5,1).. controls (-1,4) and (1,1) .. (-3,-.5);%
 \draw[b5 link] (-2,-2)--(-2,-4)--(-.7,-4);\draw[b5 link] (0,-1.5)--(0,-4)--(1.3,-4);%
\end{scope}}%
\providecommand{\BasicsFiveCat}[3]{\begin{scope}[shift={(#1,#2)},scale=#3]%
 \draw[draw=B5Gray,fill=white,line width=.8pt] (0,-1) ellipse (3.3 and 4.3);%
 \draw[draw=B5Gray,fill=white,line width=.8pt,line join=round]%
  (-3,3)--(-3,7)--(-1.2,5.8).. controls (0,6.5) and (1,6.5) .. (2,5.8)%
  --(3.3,7)--(3.3,3).. controls (2,1.3) and (-1.5,1.3) .. (-3,3);%
 \fill[B5Gray] (-1.3,4) circle (.3);\fill[B5Gray] (1.5,4) circle (.3);%
 \draw[b5 link] (0,3.4)--(0,2.7);\draw[b5 link] (-2,-2)--(-2,-5);\draw[b5 link] (2,-2)--(2,-5);%
 \draw[draw=B5Gray,line width=.8pt] (-2.5,-3).. controls (-8,-5) and (-8,-.5) .. (-5,-1);%
\end{scope}}%
\providecommand{\BasicsFiveRobot}[3]{\begin{scope}[shift={(#1,#2)},scale=#3]%
 \fill[B5Gray,rounded corners=.4mm] (-2.7,-1.6) rectangle (-1.7,1.6);%
 \fill[B5Gray,rounded corners=.4mm] (1.7,-1.6) rectangle (2.7,1.6);%
 \draw[draw=B5Gray,fill=white,line width=1pt,rounded corners=.7mm] (-1.8,-2.3) rectangle (1.8,2.3);%
 \draw[draw=B5Gray,fill=B5Yellow,line width=.8pt] (0,.6) circle (.65);%
\end{scope}}%
\BasicsFiveSetup%
\begin{tikzpicture}[b5 picture]%
 \path[use as bounding box] (0,-10) rectangle (168,81);%
 \node[b5 node,minimum width=24mm,minimum height=20mm,fill=B5Ice] at (13,45) {};%
 \BasicsFiveBird{15}{44}{1}{white}\node at (13,32) {原图};%
 \node[b5 node,minimum width=25mm,minimum height=20mm,fill=B5Ice] at (47,65) {};%
 \begin{scope}\clip[rounded corners=1.5mm] (34.5,55) rectangle (59.5,75);%
  \BasicsFiveBird{43}{62}{1.7}{white}\end{scope}%
 \node at (47,52) {裁剪};%
 \node[b5 node,minimum width=25mm,minimum height=20mm,fill=B5Ice] at (47,38) {};%
 \BasicsFiveBird{49}{37}{1}{B5Ice}\node at (47,25) {颜色变化};%
 \node[b5 node,minimum width=25mm,minimum height=20mm,fill=B5Ice] at (47,7) {};%
 \BasicsFiveCat{49}{7}{1}\node at (47,-6) {其他原图};%
 \draw[b5 flow] (25,49).. controls (30,49) and (27,65) .. (34.5,65);%
 \draw[b5 flow] (25,41).. controls (29,41) and (29,38) .. (34.5,38);%
 \draw[draw=B5Gray,fill=none,line width=1pt,rounded corners=1.5mm]%
  (81,74)--(100,59)--(119,74)--(119,11)--(100,24)--(81,11)--cycle;%
 \BasicsFiveNet{88}{42}{2.4}%
 \node at (100,79) {共享编码器 $f_{\bm\theta}$};%
 \draw[b5 flow] (59.5,65).. controls (68,65) and (72,62) .. (81,62);%
 \draw[b5 flow] (59.5,38)--(81,38);%
 \draw[b5 flow] (59.5,7).. controls (70,7) and (71,17) .. (81,17);%
 \node[b5 node,minimum width=35mm,minimum height=59mm] at (150,42) {};%
 \node at (150,65) {表示空间};%
 \foreach \yy in {62,38,17}{\draw[b5 flow] (119,\yy)--(132.5,\yy);}%
 \draw[draw=B5Blue,fill=B5Blue,line width=1pt] (142,49) circle (1.6);%
 \draw[draw=B5Blue,fill=white,line width=1pt] (151,47) circle (1.6);%
 \draw[draw=B5Coral,fill=B5Rose,line width=1pt] (161,29) rectangle (164,32);%
 \node at (142,54) {$\bm z_1$};\node at (153,52) {$\bm z_2$};%
 \node at (162,25) {$\bm z^{-}$};%
 \draw[draw=B5Blue,line width=1.1pt,{Stealth[length=1.4mm]}-{Stealth[length=1.4mm]}] (144,49)--(149,47.5);%
 \node at (147,59) {拉近};%
 \draw[draw=B5Coral,line width=1.1pt,dashed,{Stealth[length=1.4mm]}-{Stealth[length=1.4mm]}] (152,45)--(161,33);%
 \node[text=B5Coral] at (160,42) {分离};%
\end{tikzpicture}%
\endgroup%

  \caption{自监督对比学习的教学示意。同一幅鸟的原图经过保留动物可识别性的裁剪与颜色变化，形成两个视图；它们与另一幅猫图的视图经共享参数的编码器得到表示。实心点与空心圆表示同源匹配，方点表示其他原图的对照；箭头只示意同源拉近与对照相对分离的训练目标。动物图元和表示位置均为教学构造，不是实测影像或表示分布。}
  \label{fig:contrastive-learning-process}
\end{figure*}
\FloatBarrier

匹配关系还可以有别的来源。监督式对比学习用类别标签决定哪些样本应当接近
\scite{khosla2020}。图像与文字也可以配对，例如一张图片和描述它的文字。因此，“对比”说明
怎样构造表示学习目标，“自监督”说明监督信号从哪里来；SimCLR同时符合两种描述，前面的
掩码预测则属于不使用对比目标的自监督方法。

这种学习依赖于我们选定的关系。若裁剪已改变图片含义，原本的正样本对就可能不再适合匹配；
不同原图也可能拍到了语义相近的对象，把它们全部作为负样本会与后续分类需要发生冲突。
使用对比目标时，需要检查这些关系是否保留了目标任务真正关心的信息。

\subsection{主动学习}
\label{sec:active-learning}

对比学习利用已有样本之间的关系。如果还能请人补充标签，但预算有限，就需要决定先标注哪些
样本。\term{主动学习}（active learning）让学习者根据当前知识挑选需要询问的样本，再请
人或其他可靠来源给出答案。它希望新增标签能够纠正常见错误，或补足尚未覆盖的情形，使有限
预算发挥更大作用。

在常见的样本池设定中，我们有一个未标注样本池、一个当前模型、可以回答查询的标注来源，
以及有限的查询预算。选样需要这些资源的支持：若连相关输入都收集不到，或没有可靠的回答来源，
仅靠主动选样无法获得缺失的信息。

一项图像主动学习研究从20张各数字类别数量均衡的MNIST标注图片开始，每轮从未标注池选择10张图片获取标签，
再重新学习；研究比较了不同选择准则与随机选样\scite{gal2017}。基准实验中的“请求标签”是
揭示数据集中已经存在但对学习者隐藏的标签，并不表示实验现场真的有标注员实时参与。真实应用
可以由医生、工程师或普通标注者提供答案。图~\ref{fig:active-learning-loop}中的循环体现了
主动学习与一次性标注的区别：获得新标签、更新模型之后，下一轮的选样也可以随之改变。

\begin{figure*}[htbp]
  \centering
  % 两轮池内七个数字图元沿用既有示意；边界不是实测拟合。
\begingroup%
% 第1—3章局部矢量图元：单色黄、双色蓝红、三色蓝红黄。
\providecommand{\BasicsFiveSetup}{%
 \colorlet{B5Blue}{FigureBlue}\colorlet{B5Coral}{FigureRed}%
 \colorlet{B5Yellow}{FigureYellow}\definecolor{B5Gray}{HTML}{4C4D4F}%
 \colorlet{B5Ice}{FigureBlueFill}\colorlet{B5Rose}{FigureRedFill}%
 \colorlet{B5Cream}{FigureYellowFill}%
 \tikzset{b5 picture/.style={x=1mm,y=1mm,font=\figurefont\mdseries\fontsize{8.5}{10.5}\selectfont,text=B5Gray,line width=1pt},%
 b5 node/.style={draw=B5Gray,fill=white,line width=1pt,rounded corners=1.5mm,inner sep=3pt,font=\figurefont\mdseries\fontsize{8.5}{10.5}\selectfont,text=B5Gray,align=center},%
 b5 flow/.style={draw=B5Gray,line width=1.1pt,-{Stealth[length=2.2mm,width=1.4mm]}},%
 b5 link/.style={draw=B5Gray,line width=.7pt},%
 b5 feedback/.style={b5 flow,draw=B5Coral,dash pattern=on 2.5mm off 1.5mm}}}%
% 3—3—1网络仅为图标，不指定真实架构；各层用实心蓝、红、黄，连线中性灰。
\providecommand{\BasicsFiveNet}[3]{\begin{scope}[shift={(#1,#2)},scale=#3]%
 \foreach \a in {-3,0,3}{\foreach \b in {-3,0,3}{\draw[b5 link] (0,\a)--(5,\b);}}%
 \foreach \a in {-3,0,3}{\draw[b5 link] (5,\a)--(10,0);}%
 \foreach \a in {-3,0,3}{%
  \draw[draw=B5Gray,line width=.8pt,fill=B5Blue] (0,\a) circle (1.1);%
  \draw[draw=B5Gray,line width=.8pt,fill=B5Coral] (5,\a) circle (1.1);}%
 \draw[draw=B5Gray,line width=.8pt,fill=B5Yellow] (10,0) circle (1.1);%
\end{scope}}%
% 鸟与猫保留教学对象，不生成或模拟真实数据样本。
\providecommand{\BasicsFiveBird}[4]{\begin{scope}[shift={(#1,#2)},scale=#3]%
 \draw[draw=B5Gray,fill=#4,line width=.8pt,line join=round]%
  (-6,-1)--(-10,-2)--(-6,2).. controls (-3,3) and (-2,6) .. (1,6)%
  .. controls (4,6) and (5,4) .. (4,2).. controls (4,-1) and (-1,-3) .. (-6,-1);%
 \draw[draw=B5Gray,fill=white,line width=.7pt] (4,3)--(6,3.8)--(4,4.1)--cycle;%
 \fill[B5Gray] (2.7,4.3) circle (.35);%
 \draw[b5 link] (-5,1).. controls (-1,4) and (1,1) .. (-3,-.5);%
 \draw[b5 link] (-2,-2)--(-2,-4)--(-.7,-4);\draw[b5 link] (0,-1.5)--(0,-4)--(1.3,-4);%
\end{scope}}%
\providecommand{\BasicsFiveCat}[3]{\begin{scope}[shift={(#1,#2)},scale=#3]%
 \draw[draw=B5Gray,fill=white,line width=.8pt] (0,-1) ellipse (3.3 and 4.3);%
 \draw[draw=B5Gray,fill=white,line width=.8pt,line join=round]%
  (-3,3)--(-3,7)--(-1.2,5.8).. controls (0,6.5) and (1,6.5) .. (2,5.8)%
  --(3.3,7)--(3.3,3).. controls (2,1.3) and (-1.5,1.3) .. (-3,3);%
 \fill[B5Gray] (-1.3,4) circle (.3);\fill[B5Gray] (1.5,4) circle (.3);%
 \draw[b5 link] (0,3.4)--(0,2.7);\draw[b5 link] (-2,-2)--(-2,-5);\draw[b5 link] (2,-2)--(2,-5);%
 \draw[draw=B5Gray,line width=.8pt] (-2.5,-3).. controls (-8,-5) and (-8,-.5) .. (-5,-1);%
\end{scope}}%
\providecommand{\BasicsFiveRobot}[3]{\begin{scope}[shift={(#1,#2)},scale=#3]%
 \fill[B5Gray,rounded corners=.4mm] (-2.7,-1.6) rectangle (-1.7,1.6);%
 \fill[B5Gray,rounded corners=.4mm] (1.7,-1.6) rectangle (2.7,1.6);%
 \draw[draw=B5Gray,fill=white,line width=1pt,rounded corners=.7mm] (-1.8,-2.3) rectangle (1.8,2.3);%
 \draw[draw=B5Gray,fill=B5Yellow,line width=.8pt] (0,.6) circle (.65);%
\end{scope}}%
\BasicsFiveSetup%
\begin{tikzpicture}[b5 picture]%
 \path[use as bounding box] (0,-1) rectangle (168,94);%
 \foreach \base/\time in {49/t,5/{t+1}}{%
  \draw[draw=B5Gray,fill=white,line width=1pt] (0,\base) rectangle (54,\base+36);%
  \draw[b5 flow,line width=.9pt] (4,\base+4)--(51,\base+4) node[below]{$x_1$};%
  \draw[b5 flow,line width=.9pt] (4,\base+4)--(4,\base+33) node[left]{$x_2$};%
  \node at (27,\base+40) {样本池 $\time$};%
  \foreach \xx/\yy/\digit in {12/26/3,12/18/3,17/10/3,25/19/3,41/26/8,46/19/8,39/11/8}{%
   \ifnum\digit=3\def\poolfill{B5Ice}\else\def\poolfill{B5Rose}\fi%
   \draw[fill=\poolfill,draw=none] (\xx,\base+\yy) circle (2.5);%
   \node at (\xx,\base+\yy) {$\digit$};}}%
 \draw[draw=B5Blue,line width=1.2pt] (21,53).. controls (36,63) and (27,72) .. (36,83);%
 \draw[draw=B5Blue,line width=1pt,dashed] (21,9).. controls (36,19) and (27,28) .. (36,39);%
 \draw[draw=B5Blue,line width=1.2pt] (29,9).. controls (41,19) and (34,28) .. (43,39);%
 \draw[draw=B5Yellow,line width=1.1pt] (25,68) circle (3);%
 \draw[draw=B5Blue,line width=1.1pt] (25,24) circle (3);%
 \draw[draw=B5Yellow,line width=1.1pt] (46,24) circle (3);%
 \node[b5 node,fill=B5Ice,minimum width=40mm,minimum height=24mm] (annotation) at (115,71) {};%
 \node at (115,80) {标注接口};%
 \draw[draw=B5Gray,fill=white,line width=.8pt,rounded corners=.7mm] (98,61) rectangle (114,77);%
 \draw[draw=B5Gray,line width=1.2pt,line cap=round] (102,73).. controls (111,78) and (113,69) .. (107,69).. controls (115,68) and (111,61) .. (102,64);%
 \node[b5 node,minimum width=15mm,minimum height=5mm,fill=B5Ice,inner sep=1pt] at (124,73) {$y=3$};%
 \foreach \yy in {67,63}{\draw[draw=B5Gray,fill=white,line width=.7pt,rounded corners=.5mm] (117,\yy) rectangle (132,\yy+2);}%
 \draw[b5 flow] (54,68).. controls (73,76) and (81,76) .. (annotation.west);%
 \node at (77,78) {查询};%
 \draw[draw=B5Gray,fill=white,line width=1pt] (144,34)--(144,59) arc[start angle=180,end angle=360,x radius=11,y radius=2.5]--(166,34) arc[start angle=0,end angle=-180,x radius=11,y radius=2.5];%
 \draw[draw=B5Gray,fill=white,line width=1pt] (155,59) ellipse (11 and 2.5);%
 \node at (155,54) {带标签集合 $D_L$};%
 \foreach \yy/\digit in {49/3,44/3,39/8}{%
  \node[b5 node,minimum width=18mm,minimum height=4mm,inner sep=1pt] at (155,\yy) {$\digit$};}%
 \node at (155,33.5) {$\cdots$};%
 \draw[b5 flow] (annotation.east).. controls (151,72) and (155,65) .. (155,61.5);%
 \node[b5 node,minimum width=34mm,minimum height=29mm] (retrain) at (114,17) {};%
 \node at (114,27) {重训 $\bm\theta^{(t+1)}$};\BasicsFiveNet{106}{13.5}{1.6}%
 \draw[b5 flow] (155,31.5).. controls (155,21) and (144,17) .. (retrain.east);%
 \draw[b5 flow,rounded corners=1.5mm] (retrain.west)--(70,17)--(70,24)--(54,24);%
 \draw[b5 feedback] (54,24).. controls (84,25) and (76,57) .. (annotation.south west);%
 \node[text=B5Coral] at (72,46) {下一轮};%
\end{tikzpicture}%
\endgroup%

  \caption{池式主动学习的教学示意。蓝色曲线表示当前模型的决策边界，暖色圆圈标出本轮选择的不确定样本。标注者返回标签后，样本并入带标签集合，用于重训模型。下方的点线保留更新前的边界，实线表示更新后的边界；中央带蓝色圈的数字3已获标签，右侧带暖色圈的数字8是下一轮查询对象。查询预算限制循环。数字图元、二维位置和边界均为教学构造，不表示实测分布或选样收益。}
  \label{fig:active-learning-loop}
\end{figure*}
\FloatBarrier

选择哪些样本请求标注，需要考虑这些标签能否帮助模型改进。例如，模型可能对一张严重模糊的
数字图片很不确定，但即使获得了它的标签，也未必能学到适用于其他图片的规律；若连续选中
许多写法几乎相同的数字，新增标签提供的信息也可能重复。因此，选样时除了看模型有多不确定，
还需要考虑样本是否覆盖不同情形，以及标注者能否可靠地给出答案、需要付出多少成本。

主动学习还可以与半监督学习配合。在手写数字识别中，学习者可以挑选部分图片请求真实标签，
同时通过伪标签或一致性约束利用其余未标注图片。前者决定把有限的标注预算用在哪里，后者
让没有获得人工标签的图片也能参与学习。

这个查询循环也容易让人联想到强化学习。二者确实都包含选择与反馈，但本例每次询问得到的是
所选图片的标签，评价目标是在给定标注预算下提高识别能力；强化学习则用奖励评价行动后果，
关注累计回报。我们可以进一步把多轮选样写成序贯决策问题，用强化学习寻找查询策略，
但是否这样实现，是主动学习设定之外的另一项选择。

\subsection{标签利用方式的比较}
\label{sec:efficient-knowledge-use-comparison}

同样面对标签有限的分类任务，这几种范式利用的资源不同。假设已经收集了一批花卉照片，只有
少量照片带有类别标签：半监督学习让已有标签与其余无标签照片共同服务于花卉分类；自监督
学习则从照片自身构造训练目标，先学习表示，之后仍需用类别标签或其他信息说明要识别什么。
前者直接围绕目标任务利用两类数据，后者着重解决没有额外任务标签时怎样获得学习信号。

若专家只能提供粗略判断，或已有规则能够批量产生不完全可靠的标签，弱监督学习研究怎样利用
这些信号并处理其中的错误。若专家可以准确标注部分照片，主动学习则决定优先询问哪些照片。
两者分别改变可接受的监督质量和新增标签的获取方式，不能只凭“都节省标注”就混为一谈。

对比学习关注的是怎样使用样本关系：同一照片的不同视图可以提供自监督匹配关系，类别标签
也可以决定哪些照片应当接近。因此，自监督说明目标的来源，对比说明目标的构造方式，两者
可以重叠；半监督也属于广义弱监督中的监督不完整情形。
表~\ref{tab:efficient-knowledge-use-comparison}汇总了这些差别，便于根据可用资源判断哪些
办法适用，而不是为一个系统挑选唯一名称。

\begin{table*}[htbp]
  \centering
  \small
  \begin{tabularx}{\linewidth}{@{}lXXX@{}}
    \toprule
    \textbf{范式或方法} & \textbf{可以利用的资源} & \textbf{减少标签依赖的方式} & \textbf{需要检查的条件} \\
    \midrule
    半监督学习 & 少量标签与无标签数据 & 让无标签样本也参与目标任务学习 & 数据结构与目标类别是否相符 \\
    自监督学习 & 原始内容及其关系 & 从数据中构造训练目标 & 所学规律是否对实际任务有用 \\
    弱监督学习 & 粗标签、噪声标签或规则 & 接受成本较低的不完善监督 & 标签误差及不同来源的依赖 \\
    对比学习 & 视图、配对关系或类别标签 & 用匹配关系学习可复用表示 & 匹配依据是否可靠；也可使用标签 \\
    主动学习 & 可查询的样本与标注来源 & 优先为有价值的样本获取标签 & 查询收益、样本覆盖与标注成本 \\
    \bottomrule
  \end{tabularx}
  \caption{围绕有限标注比较知识的高效利用。减少人工任务标签并不等于减少全部训练成本，
  也不保证这些方法在任意数据上都能提高效果。}
  \label{tab:efficient-knowledge-use-comparison}
\end{table*}
\FloatBarrier

\section{知识的演进与迁移}
\label{sec:knowledge-evolution-and-transfer-overview}

标签预算并不是经验利用的全部问题。模型可能先在白天拍摄的道路图像上学习，之后要处理夜间
图像；也可能原先只判断物体类别，后来还要估计距离。已经掌握的知识是否仍然适用，取决于
数据和任务分别发生了什么变化，以及新旧能力之间需要保持什么关系。

这里要区分两种变化。\emph{数据变化}不只是换了几张照片，还可能是拍摄条件、样本组成或
输入与目标之间的统计关系改变；从同一分布抽取一批新样本，不等于分布已经改变。
\emph{任务变化}则涉及需要输出什么、区分哪些类别或按什么目标评价，例如从物体识别转为
距离估计。相同图像可以服务不同任务，同一任务也可以面对不同的数据分布。

持续学习、多任务学习、迁移学习和元学习从不同角度处理知识的演进与迁移：经验可以沿时间
积累，可以在多个任务间共享，可以从来源问题用于目标问题，也可以用来改进后续学习过程。
比较时，除了数据与任务是否变化，还要看任务怎样组织、历史数据能否访问，以及最终评价哪些
能力。下面从长期积累知识的情形展开。

\subsection{持续学习}
\label{sec:continual-learning}

长期运行的系统往往无法一次获得未来所需的全部经验。用户兴趣、拍摄条件会变化，新类别也
会不断出现。系统需要适应这些变化，同时保留仍有用的旧能力，这就是\term{持续学习}
（continual learning）。

新旧能力可能依靠同一组参数。只用新数据更新时，旧任务依赖的参数也会改变，旧类别上的表现
可能明显下降，这称为\term{灾难性遗忘}（catastrophic forgetting）。因此，学会新内容后，
还必须回头检查旧能力。

以图像识别为例，可以让模型分批接收新类别，并要求它在每个阶段识别截至当时见过的全部
类别，预测时不告知图片属于哪个阶段。这种设定称为\term{类别增量学习}
（class-incremental learning）。iCaRL是处理这一
问题的方法，其实验在CIFAR-100和ImageNet数据上逐步加入类别，同时
限制能够保存的旧样本数量\scite{rebuffi2017}。图~\ref{fig:continual-learning-timeline}展示了
这一过程：每个阶段不只考察新类别是否学会，还检查旧类别是否被遗忘。

\begin{figure}[!htb]
  \centering
  \bookdraftgraphic{figures/scenes/v1r5-continual.png}
  \caption{持续学习的类别增量过程。不同类别的图像按阶段到来；每个阶段都要学习新增类别，并在
  截至当前出现过的全部类别上检查能力。三个模型图标沿用相同轮廓与分区，内部逐步保留的填充示意知识积累；分区不对应真实神经元或精确记忆容量。向过去回指的箭头表示对旧知识的保留要求。}
  \label{fig:continual-learning-timeline}
\end{figure}
\FloatBarrier

讨论持续学习时，需要说明过去的数据还能保留多少、系统是否会被告知一个新任务已经开始，
以及何时检查新旧能力。如果每个阶段都能使用全部历史数据，就可以通过重新训练复习旧内容；
如果只能保留少量旧样本，维持旧能力就受到更强的限制。持续学习也不限于新增类别，还可以
处理传感器变化、用户兴趣变化或依次出现的新任务。

任务信息是否可见，还会改变问题难度。若陆续学习不同任务，预测时明确告知当前样本应由哪个
任务处理，就得到\term{任务增量学习}（task-incremental learning）；系统可以据此选择相应
的输出。若识别目标与输出含义保持不变，数据分布陆续变化，而且预测时不提供任务身份，
则是典型的\term{领域增量学习}（domain-incremental learning）。它与类别增量、任务增量
构成持续学习中三种常用的比较设定\scite{vandeven2019}。

例如，白天与夜间图像都要求识别同一组交通标志，可以作为领域增量问题；逐步加入原先没有的
标志类别，并要求在全部已见类别中作出判断，则是类别增量问题。如果每次预测都先告知属于
哪一组识别任务，系统面对的就更接近任务增量设定。三者都需要按要求保留旧能力，但可用提示
和预测范围不同，实验结果不能脱离这些条件比较。若夜间识别改善，却完全失去白天识别能力，
就没有满足同时应对两种环境的长期目标。

\subsection{多任务学习}
\label{sec:multi-task-learning}

持续学习受到经验到达顺序与历史数据可访问性的限制。若多个相关任务的数据已经可用，
则可以在训练期间共同利用它们。例如，分析街景既要知道哪里是车辆和道路，也要估计它们
有多远，两项任务都可能用到
物体边界等信息。\term{多任务学习}（multi-task learning）让相关任务共享表示或其他知识，
共同学习各自的目标。这样，一个任务的反馈可能帮助另一个任务发现需要的特征，减少各任务
分别学习相同规律的负担\scite{caruana1997}。这种帮助来自学习过程中的信息共享；只是同时
运行多个互不联系的模型，并不能产生这种作用。

具体到街景，逐像素判断道路、车辆或建筑等类别，称为\term{语义分割}（semantic segmentation）；
估计各像素对应物体的距离，称为\term{深度估计}（depth estimation）。Kendall等人还将
区分不同物体个体的\term{实例分割}（instance segmentation）纳入联合学习，并研究怎样
平衡各任务反馈的影响\scite{kendall2018}。

图~\ref{fig:multi-task-learning-process}只展示前两项任务：同一图像经过共享的处理部分，
再分别产生语义区域和深度预测。两项预测各自与相应标注比较，反馈又共同影响共享部分。
因此，一个任务学到的边界信息有机会用于另一个任务，而每项任务仍保有自己的预测要求。

多任务学习不要求任务随时间发生变化。上面的系统可以始终只学习语义分割与深度估计这两个
任务，数据也可以预先收集完毕；需要组织的是多个目标之间的知识共享。两个任务可以使用同一批
输入及不同标注，也可以拥有各自的数据集，关键在于一个任务的学习能够影响另一个任务。

这与持续学习的区别不能只看“是否有多个任务”。多任务学习通常能够在训练期间反复利用各
任务的数据，即使每次更新交替选择一个任务，也仍然可以是联合学习。持续学习则按先后到来的
经验更新，旧数据常常无法完整访问，因此还要专门检查旧能力是否被新学习破坏。

\begin{figure}[!htb]
  \centering
  \bookdraftgraphic{figures/scenes/v1r5-multi-task-learning.png}
  \caption{多任务学习的过程示意。同一街景及不同任务的标注共同影响共享表示，系统分别输出
  语义区域和远近关系；两个结果都是概念示意，不是论文实验的实际输出。}
  \label{fig:multi-task-learning-process}
\end{figure}
\FloatBarrier

相关任务可能互相帮助，也可能发生干扰：改善一个任务的参数调整可能损害另一个任务。
某一任务数据更多、目标更容易，还可能主导共享表示，使其他任务受损。因此，任务越多并不
一定越好，需要分别检查各任务的效果，不能只看一个汇总指标。

\subsection{迁移学习}
\label{sec:transfer-learning}

共同训练多个任务是一种利用相关性的方式。如果已有模型学过相关问题，现在希望借助这些
知识解决一个目标问题，就可以把来源与目标区分开。\term{迁移学习}（transfer learning）把已有任务或领域中学到的表示、参数等知识
用于目标问题，减少需要从头学习的内容。它既可用于标签有限的新任务，也可用于加快学习或
适应数据分布变化；关键是来源知识中哪些部分对目标问题仍然有用。

例如，一个模型先从ImageNet中的大量通用图片学习图像特征，再利用较少的带标签花卉照片学习
区分花卉类别。来源阶段提供边缘、纹理和形状等可复用表示，目标阶段仍需要用花卉任务自己的
标签确定哪些差异重要。图~\ref{fig:transfer-learning-process}中，从来源任务通向花卉任务的
箭头表示复用已经学到的表示或模型参数：花卉分类从这个基础继续学习，而不必从完全没有图像
经验的状态开始。图中的花卉任务是用于说明迁移过程的教学情景。

\begin{figure}[!htb]
  \centering
  \bookdraftgraphic{figures/scenes/v1r3-transfer-learning.png}
  \caption{迁移学习的过程示意。通用图像任务形成可复用的表示，目标任务再用少量花卉标签调整
  这些知识，最后识别新的花卉照片；来源知识与目标数据共同决定结果。}
  \label{fig:transfer-learning-process}
\end{figure}
\FloatBarrier

花卉分类仍需要自己的标签，是因为来源任务不会自动告诉模型哪些差异值得关注。边缘、颜色
和纹理可能是识别一般物体与花卉的共同基础，但足以区分汽车和动物的特征，未必足以区分两种
外形相近的花。目标标签帮助模型进一步抓住这些细微差异。

图像神经网络的多层处理也体现了这种差别：较早的层常提取边缘、颜色和局部纹理，较后的层
再把这些信息组合成与训练任务有关的特征。不同层与来源任务的联系不同，可直接迁移的程度
也不同。针对图像网络的研究观察到，较高层的特征通常更依赖原任务，来源与目标越不相似，迁移越受限制
\scite{yosinski2014}。这并不意味着高层特征一概不能复用，而是需要结合目标任务决定保留和
调整哪些部分。

把已有模型的参数作为起点，再用目标任务数据继续调整，称为\term{微调}（fine-tuning）。
另一种方式是固定已有表示，只学习目标任务所需的输出部分。它们是复用知识的具体办法；
选择哪一种，还要考虑目标数据量、来源与目标的差异以及允许调整的范围。

与主动学习相比，迁移学习主要增加的是目标学习者可以利用的先前知识，不必向外界请求新的
目标标签。两者可以配合：先迁移通用图像表示，再挑选当前最需要确认的花卉照片请求标注。
这种组合能否奏效，既取决于来源知识是否相关，也取决于新增标签是否补足了目标任务的缺口。

迁移不要求任务与数据分布同时变化。前面的通用物体分类转向花卉分类，改变了需要区分的
类别；白天训练的物体识别模型用于夜间图像时，类别可以保持不变，改变的主要是输入分布。
比较迁移设定时，应说明来源与目标究竟在哪些方面不同，而不是只说模型换了一批数据。

在街景例子中，若先用深度估计学到视觉表示，再将它用于语义分割，最终关注语义分割的改善，
就突出了来源任务对目标任务的帮助。若两项监督共同影响共享表示，并分别评价两项任务，
则突出了多任务学习。两者可以重叠，迁移也可以借助联合训练实现，不能仅凭训练是否先后判断。
与持续学习相比，一次迁移可以只要求目标任务表现良好，不必继续保留来源任务的全部能力。

若复用来源知识反而使目标任务学得更差，就发生了\term{负迁移}（negative transfer）。判断
迁移效果，需要与不采用这些知识时的目标任务表现比较，不能只看来源模型有多大或用了多少数据。

任务不变而数据来源变化时，还需要明确目标数据何时可用。在道路识别中，白天、雨天与夜间
的图像可能呈现不同的统计规律。这里用\term{域}（domain，也称领域）描述输入空间及其
数据分布，用\term{源域}（source domain）指提供已有知识的数据来源，用\term{目标域}
（target domain）指最终希望模型有效工作的环境。域不能只按数据集名称区分，同一数据集中
也可能包含多个环境。

假设模型已从白天道路图像学到识别能力，现在能提前收集一批夜间图像，就可以利用这些目标域
数据调整模型。\term{领域自适应}（domain adaptation，也称领域适应）关注借助可获得的
目标域数据应对源域与目标域的分布差异。若目标域提供标签，就是有监督领域自适应；若源域有
标签、目标域只提供无标签输入，则是\term{无监督领域自适应}（unsupervised domain
adaptation）。域对抗方法尝试学习较难据以区分数据来自哪个域的表示，同时保留源域分类能力，
以帮助迁移\scite{ganin2016}。这里的“无监督”限定目标域的标签条件，并不表示整个学习过程
没有任何标签；缩小表示差异也不自动保证目标识别正确。

但部署地点可能尚未确定，训练时根本收集不到目标域数据。\term{领域泛化}
（domain generalization）要求仅利用一个或多个源域学习，使模型能够在未见目标域上工作
\scite{muandet2013}。例如，只用若干城市的白天和雨天图像准备识别模型，之后直接用于此前
没有接触过的夜间道路。这时不能提前根据夜间图像调整表示，也不能假定训练已经覆盖所有
将来会遇到的变化。与从同一分布抽取新样本的泛化评价相比，领域泛化还要求应对训练中未见的
环境差异。

这一边界同样约束模型选择。本章采用的领域泛化设定不允许用目标域数据选择模型、调整超参数
或决定何时停止训练；这些决定需要依靠源域数据，例如留出一个源域来模拟环境变化。
Gulrajani与Lopez-Paz的研究专门比较了模型选择方式对领域泛化评价的影响
\scite{gulrajani2021}。目标域数据在这里用于最终检验。模型当然需要读取待预测图像才能输出
结果，但读取输入进行预测与利用它更新模型是两回事。

领域泛化可以利用不同源域之间较稳定的特征，例如尝试保留物体形状，减少对某一背景的依赖。
但“在源域间稳定”不等于“在所有未来环境中都稳定”，源域数量增加也不必然带来目标域改善。
若新环境中的输入或标签规律任意改变，源域经验不足以保证识别成功。领域泛化仍需要来源与
目标之间存在可利用的联系，其可学习条件将在第~\ref{sec:domain-generalization}节讨论。

部署之后又多了一种可能：虽然训练时见不到夜间图像，系统使用时却能接触这些输入。
\term{测试时适应}（test-time adaptation）允许在测试或部署阶段利用这些数据调整模型参数
或统计量。常见设定不提供目标标签；例如，Tent利用测试输入及模型预测更新部分参数，
以降低预测的不确定性\scite{wang2021tent}。不过，更有把握不一定意味着更正确，错误预测
也可能被更新过程强化。需要说明允许使用哪些测试输入、何时更新，以及不同样本之间是否
保留更新结果。

表~\ref{tab:domain-transfer-settings}以目标域信息何时可用比较这些设定。若训练阶段按领域
泛化准备模型，部署时又继续适应，就同时涉及两个阶段的要求；评价时应明确测试时发生了
更新，不能把适应后的结果当成固定模型直接泛化的结果。

\begin{table*}[htbp]
  \centering
  \small
  \begin{tabularx}{\linewidth}{@{}lXX@{}}
    \toprule
    \textbf{设定} & \textbf{可用的目标域信息} & \textbf{关键条件} \\
    \midrule
    有监督领域自适应 & 适应时已有目标域带标签样本 & 用目标监督调整模型，另留样本评价 \\
    无监督领域自适应 & 适应时已有目标域无标签输入 & 源域标签与目标域输入共同参与学习 \\
    领域泛化 & 训练及模型选择时均无目标域数据 & 本章标准设定在目标域直接预测，不更新模型 \\
    测试时适应 & 测试或部署时利用目标输入；常见设定无标签 & 按约定更新参数或统计量，并说明数据与时序限制 \\
    \bottomrule
  \end{tabularx}
  \caption{目标域数据的可用性区分领域自适应、领域泛化与测试时适应。前两行区分目标监督条件，
  最后一行强调适应发生的阶段；它们不是完全互斥的分类。}
  \label{tab:domain-transfer-settings}
\end{table*}
\FloatBarrier

迁移到目标任务时，还需要说明目标一侧能够提供多少带标签示例。前面的花卉分类可以使用较多
标注，也可以只给每个新类别几个示例；如果连这几个示例都没有，就需要类别描述等其他信息
来表达目标。在这里的分类情景中，\term{少样本学习}（few-shot learning）指利用极少的
带标签示例识别新类别；\term{零样本学习}（zero-shot learning）则不提供目标类别的带标签
训练示例，而借助先前知识及类别描述等辅助信息。两个名称限定的是目标侧的数据条件，
并没有规定知识必须怎样获得或迁移。

类别描述怎样变成识别依据？模型需要先学会图像内容与文字含义之间的对应关系。
CLIP是一种同时学习图像和文本表示的方法\scite{radford2021}。训练数据由图片及其配套描述组成，学习过程让
相互匹配的图像与文字表示更接近，并与同批次中不匹配的图文组合区分开。这沿用了
第~\ref{sec:contrastive-learning}节的对比思路，只是匹配双方换成了图像与文字。模型积累的
不仅是图像特征，还包括哪些文字能够描述哪些视觉内容。

使用时，可以把候选类别写成简短描述，例如“a photo of a cat”（一张猫的照片）和
“a photo of a dog”（一张狗的照片），再比较新图像与各段文字的匹配程度，以最匹配的描述
确定类别。这个过程不需要再提供该目标分类任务的带标签训练图片，因此是一种零样本分类方式。
这里的“零”不表示系统从未学习过任何数据：识别依赖先前的图文学习，也不能保证预训练语料
没有包含相关概念。

零样本学习与领域泛化也不能互相替代。前者在这里限制目标类别的带标签训练示例，后者限制
目标域数据在训练及模型选择中的使用。一个模型可能对新类别进行零样本分类，却已见过同一
拍摄环境的其他图像；也可能在新环境中做领域泛化，而要识别的类别在源域中早已有标签。

\subsection{元学习}
\label{sec:meta-learning}

假设系统经常需要识别新的字符，而且每次只能看到几个带标签示例。迁移已有表示可以提供
帮助，但我们还可以提前练习这个适应过程：让模型反复面对不同的小任务，检查它从少量示例
中学得怎样，再用这些结果改进下一次学习。

\term{元学习}（meta-learning）利用一组任务上的学习经历，改善后续学习的方式或效果。
少样本适应是它的一种用途，更广义的目标还可以包括学习速度或稳健性\scite{hospedales2022}。
它常被概括为“学习如何学习”，意思是从跨任务经验中改进学习过程；具体系统仍然使用预先
规定的模型、目标和更新方式。

要实现这一目标，一种做法是让模型在许多训练任务上反复经历少样本适应，再根据适应后的效果
改善共同的学习起点。这里的“起点”是模型开始适应某个任务时使用的一组参数。
\term{模型无关元学习}（model-agnostic meta-learning, MAML）是这种思路的代表：它根据模型
从共同起点进行少量调整后，在相应任务的新样本上表现得怎样，来改进这组起始参数\scite{finn2017}。

MAML的一组实验使用Omniglot手写字符数据集，每次选少数几个字符类别，组成一项识别任务。
各任务分别从共同起点出发，用少量带标签字符调整参数，再用同一任务的另一批样本检查能否
识别新的写法。这些适应后的检查结果汇合起来，用于改进共同起点；改进后，再继续练习其他
任务。这种跨任务改进学习能力的训练阶段称为\term{元训练}（meta-training）。

图~\ref{fig:meta-learning-process}中要区分两种检查。上方各任务的检查样本虽然没有直接
用于该任务的少量适应，却参与了共同起点的学习，仍属于元训练数据。下方的最终评价使用
未参加元训练的新任务；在上述字符识别实验中，这意味着新的字符类别。评价考察模型能否
从少量新示例中学会新任务。

\begin{figure*}[!t]
  \centering
  % 训练检查参与起点更新；新任务检查只评分，两种检查不混同。
% 样本卡片先竖直落到目标中心线，共用水平干线进入同一个西侧端点。
\begingroup%
% 第1—3章局部矢量图元：单色黄、双色蓝红、三色蓝红黄。
\providecommand{\BasicsFiveSetup}{%
 \colorlet{B5Blue}{FigureBlue}\colorlet{B5Coral}{FigureRed}%
 \colorlet{B5Yellow}{FigureYellow}\definecolor{B5Gray}{HTML}{4C4D4F}%
 \colorlet{B5Ice}{FigureBlueFill}\colorlet{B5Rose}{FigureRedFill}%
 \colorlet{B5Cream}{FigureYellowFill}%
 \tikzset{b5 picture/.style={x=1mm,y=1mm,font=\figurefont\mdseries\fontsize{8.5}{10.5}\selectfont,text=B5Gray,line width=1pt},%
 b5 node/.style={draw=B5Gray,fill=white,line width=1pt,rounded corners=1.5mm,inner sep=3pt,font=\figurefont\mdseries\fontsize{8.5}{10.5}\selectfont,text=B5Gray,align=center},%
 b5 flow/.style={draw=B5Gray,line width=1.1pt,-{Stealth[length=2.2mm,width=1.4mm]}},%
 b5 link/.style={draw=B5Gray,line width=.7pt},%
 b5 feedback/.style={b5 flow,draw=B5Coral,dash pattern=on 2.5mm off 1.5mm}}}%
% 3—3—1网络仅为图标，不指定真实架构；各层用实心蓝、红、黄，连线中性灰。
\providecommand{\BasicsFiveNet}[3]{\begin{scope}[shift={(#1,#2)},scale=#3]%
 \foreach \a in {-3,0,3}{\foreach \b in {-3,0,3}{\draw[b5 link] (0,\a)--(5,\b);}}%
 \foreach \a in {-3,0,3}{\draw[b5 link] (5,\a)--(10,0);}%
 \foreach \a in {-3,0,3}{%
  \draw[draw=B5Gray,line width=.8pt,fill=B5Blue] (0,\a) circle (1.1);%
  \draw[draw=B5Gray,line width=.8pt,fill=B5Coral] (5,\a) circle (1.1);}%
 \draw[draw=B5Gray,line width=.8pt,fill=B5Yellow] (10,0) circle (1.1);%
\end{scope}}%
% 鸟与猫保留教学对象，不生成或模拟真实数据样本。
\providecommand{\BasicsFiveBird}[4]{\begin{scope}[shift={(#1,#2)},scale=#3]%
 \draw[draw=B5Gray,fill=#4,line width=.8pt,line join=round]%
  (-6,-1)--(-10,-2)--(-6,2).. controls (-3,3) and (-2,6) .. (1,6)%
  .. controls (4,6) and (5,4) .. (4,2).. controls (4,-1) and (-1,-3) .. (-6,-1);%
 \draw[draw=B5Gray,fill=white,line width=.7pt] (4,3)--(6,3.8)--(4,4.1)--cycle;%
 \fill[B5Gray] (2.7,4.3) circle (.35);%
 \draw[b5 link] (-5,1).. controls (-1,4) and (1,1) .. (-3,-.5);%
 \draw[b5 link] (-2,-2)--(-2,-4)--(-.7,-4);\draw[b5 link] (0,-1.5)--(0,-4)--(1.3,-4);%
\end{scope}}%
\providecommand{\BasicsFiveCat}[3]{\begin{scope}[shift={(#1,#2)},scale=#3]%
 \draw[draw=B5Gray,fill=white,line width=.8pt] (0,-1) ellipse (3.3 and 4.3);%
 \draw[draw=B5Gray,fill=white,line width=.8pt,line join=round]%
  (-3,3)--(-3,7)--(-1.2,5.8).. controls (0,6.5) and (1,6.5) .. (2,5.8)%
  --(3.3,7)--(3.3,3).. controls (2,1.3) and (-1.5,1.3) .. (-3,3);%
 \fill[B5Gray] (-1.3,4) circle (.3);\fill[B5Gray] (1.5,4) circle (.3);%
 \draw[b5 link] (0,3.4)--(0,2.7);\draw[b5 link] (-2,-2)--(-2,-5);\draw[b5 link] (2,-2)--(2,-5);%
 \draw[draw=B5Gray,line width=.8pt] (-2.5,-3).. controls (-8,-5) and (-8,-.5) .. (-5,-1);%
\end{scope}}%
\providecommand{\BasicsFiveRobot}[3]{\begin{scope}[shift={(#1,#2)},scale=#3]%
 \fill[B5Gray,rounded corners=.4mm] (-2.7,-1.6) rectangle (-1.7,1.6);%
 \fill[B5Gray,rounded corners=.4mm] (1.7,-1.6) rectangle (2.7,1.6);%
 \draw[draw=B5Gray,fill=white,line width=1pt,rounded corners=.7mm] (-1.8,-2.3) rectangle (1.8,2.3);%
 \draw[draw=B5Gray,fill=B5Yellow,line width=.8pt] (0,.6) circle (.65);%
\end{scope}}%
\BasicsFiveSetup%
\begin{tikzpicture}[b5 picture]%
 \path[use as bounding box] (0,-18) rectangle (168,97);%
 \node[anchor=west] at (0,90) {元训练};%
 \node[b5 node,minimum width=26mm,minimum height=12mm,fill=B5Cream] (start) at (16,55) {共同起点 $\bm\theta$};%
 \foreach \yy/\ii/\aa/\bb in {75/1/猫/狗,55/2/车/船,35/3/鸟/鱼}{%
  \node at (52.5,\yy+17) {支持样本};%
  \node[b5 node,minimum width=11mm,minimum height=8mm,fill=B5Ice] at (46,\yy+10) {\aa};%
  \node[b5 node,minimum width=11mm,minimum height=8mm,fill=B5Ice] at (59,\yy+10) {\bb};%
  \foreach \sx in {46,59}{\draw[b5 link] (\sx,\yy+6)--(\sx,\yy)--(66,\yy);}%
  \node[b5 node,minimum width=36mm,minimum height=18mm] (task\ii) at (91,\yy) {};\BasicsFiveNet{77}{\yy}{1.65}%
  \node at (102.4,\yy) {$\bm\theta'_{\tau_{\ii}}$};%
  \draw[b5 link,rounded corners=1.5mm] (start.east)--(34,55)--(34,\yy)--(66,\yy);%
  \draw[b5 flow] (66,\yy)--(task\ii.west);%
  \node at (124.5,\yy+17) {检查样本};%
  \node[b5 node,minimum width=11mm,minimum height=8mm,fill=B5Ice] at (118,\yy+10) {\aa};%
  \node[b5 node,minimum width=11mm,minimum height=8mm,fill=B5Ice] at (131,\yy+10) {\bb};%
  \foreach \qx in {118,131}{\draw[b5 link] (\qx,\yy+6)--(\qx,\yy)--(140,\yy);}%
  \node[draw=B5Gray,fill=B5Rose,circle,line width=1pt,minimum size=11mm,inner sep=0pt] (loss\ii) at (155,\yy) {$L_{\tau_{\ii}}$};%
  \draw[b5 link] (task\ii.east)--(140,\yy);%
  \draw[b5 flow] (140,\yy)--(loss\ii.west);%
  \draw[draw=B5Coral,line width=1.1pt,dashed] (loss\ii.east)--(165,\yy);}%
 \draw[b5 feedback,rounded corners=2mm] (165,75)--(165,24)--(16,24)--(start.south);%
 \node[text=B5Coral,fill=white,inner sep=1pt] at (88,24) {汇总损失，更新起点};%
 \draw[B5Gray,dashed,line width=.7pt] (0,18)--(168,18);%
 \node[anchor=west] at (0,12) {元测试};%
 \node[b5 node,minimum width=26mm,minimum height=12mm,fill=B5Cream] (learned) at (16,-6) {已学起点 $\bm\theta$};%
 \node at (52.5,13) {支持样本};%
 \node[b5 node,minimum width=11mm,minimum height=8mm,fill=B5Ice] at (46,6) {树};%
 \node[b5 node,minimum width=11mm,minimum height=8mm,fill=B5Ice] at (59,6) {山};%
 \foreach \sx in {46,59}{\draw[b5 link] (\sx,2)--(\sx,-6)--(66,-6);}%
 \node[b5 node,minimum width=36mm,minimum height=18mm] (newtask) at (91,-6) {};\BasicsFiveNet{77}{-6}{1.65}%
 \node at (102.4,-6) {$\bm\theta'_{\tau^*}$};%
 \draw[b5 link] (learned.east)--(66,-6);\draw[b5 flow] (66,-6)--(newtask.west);%
 \node at (124.5,13) {检查样本};%
 \node[b5 node,minimum width=11mm,minimum height=8mm,fill=B5Ice] at (118,6) {树};%
 \node[b5 node,minimum width=11mm,minimum height=8mm,fill=B5Ice] at (131,6) {山};%
 \foreach \qx in {118,131}{\draw[b5 link] (\qx,2)--(\qx,-6)--(140,-6);}%
 \node[draw=B5Gray,fill=B5Ice,circle,line width=1pt,minimum size=11mm,inner sep=0pt] (score) at (155,-6) {评分};%
 \draw[b5 link] (newtask.east)--(140,-6);\draw[b5 flow] (140,-6)--(score.west);%
\end{tikzpicture}%
\endgroup%

  \caption{任务内适应与跨任务更新。元训练任务从共同起点$\bm\theta$出发，用支持样本得到适应参数$\bm\theta'_{\tau}$，再汇总检查损失$L_{\tau}$更新起点；两批样本都参与元训练。元测试的新任务只用支持样本适应，检查评分不回传共同起点。任务与类别名为教学构造，同名卡片代表同一类别的不同记录，不是Omniglot原始样本。}
  \label{fig:meta-learning-process}
\end{figure*}

迁移学习也能把已有参数带到新任务，为什么还要区分元学习？可以继续用字符识别来比较。
普通的监督预训练把许多旧字符类别合在一起训练，反馈来自对这些类别的识别误差；学好后，
再把表示或参数迁移给新的字符任务。MAML则在准备阶段就把数据
组织成许多少样本任务，每次先适应、再检查，反馈来自\emph{适应之后}在该任务新样本上的
表现。两个起点即使都使用了大量旧字符数据，受到的训练要求也不同：一个以识别来源类别为
目标，另一个直接以少量示例适应后的效果为目标。

在最终使用时，两者都可能表现为“载入已有参数，再用少量新数据调整”，所以需要追溯准备
阶段的目标，才能看清差别。迁移学习关注来源知识如何帮助目标任务，元学习关注怎样用以往
任务改善学习过程；上面只是用普通监督预训练与MAML说明这种差别。两者也有交集，元学习
形成的起点用于新任务，同样可以视为知识迁移。

少样本问题也不一定要用元学习解决，普通迁移学习等方法同样可以用于这种数据条件
\scite{wang2020}。而元学习能否帮助未来任务，要看练习任务与未来任务是否相关；差异过大时，
练习形成的适应方式可能失效。新任务只用几个示例，并不表示整个系统只学过几条数据，准备
阶段仍可能需要许多历史任务。

少量示例还可以有另一种用法：把几条评论及其正面或负面标签写进语言模型的输入，再接上一条
待判断的新评论。模型据此生成答案，在这次使用中不更新参数。这种利用输入上下文中的任务
说明或示例作答的方式，称为\term{上下文学习}（in-context learning）。GPT-3的少样本
实验采用了这样的设定\scite{brown2020}。因此，看到“模型用了几个新示例”，还要检查示例
是参与参数更新，还是只进入本次输入；仅凭模型能利用输入示例，也无法判断它是否接受过以
任务适应为目标的元训练。

元学习与多任务学习都可以使用多个训练任务，但这些任务承担的作用不同。街景系统的多任务
学习要做好已有的语义分割与深度估计；上述少样本元学习则利用历史任务练习适应，并在未参与
元训练的新任务上评价学习效果。不能仅凭“训练时使用了多个任务”就认定采用了元学习。

\FloatBarrier
\subsection{数据与任务变化下的比较}
\label{sec:knowledge-evolution-comparison}

表~\ref{tab:knowledge-evolution-comparison}汇总了这四种设定。数据与任务是否变化帮助我们
定位问题，变化发生的时间、知识服务的对象以及评价要求则进一步区分学习范式。例如，多任务
学习可以面对固定的一组任务；持续学习可以只面对同一任务的数据分布变化；迁移学习与元学习
都可能用于新任务，后者还要追溯过去的训练是否在改进学习过程。

\begin{table*}[htbp]
  \centering
  \small
  \begin{tabularx}{\linewidth}{@{}lXXX@{}}
    \toprule
    \textbf{范式} & \textbf{数据怎样变化} & \textbf{任务怎样组织} & \textbf{主要评价要求} \\
    \midrule
    持续学习 & 经验顺序到来，分布可能变化；旧数据访问常受限 & 任务可以不变，也可增加类别或依次出现新任务 & 学习新内容，同时保留所需旧能力 \\
    多任务学习 & 各任务可共享输入，也可使用不同数据集；不要求随时间变化 & 一组相关任务共同学习，任务集合可以固定 & 分别检查各任务表现与共享带来的影响 \\
    迁移学习 & 来源与目标的分布可以不同 & 区分来源与目标；目标任务可相同或不同 & 来源知识是否帮助目标任务 \\
    元学习 & 用训练任务积累学习经验，再以新任务数据考察适应 & 跨任务改进学习过程；本节以少样本适应为例 & 在给定适应数据和预算下学得怎样 \\
    \bottomrule
  \end{tabularx}
  \caption{从数据、任务及评价要求比较知识的演进与迁移。各行描述本节的典型设定，
  不是由“变或不变”唯一确定的互斥类别。}
  \label{tab:knowledge-evolution-comparison}
\end{table*}
\FloatBarrier

仍以道路图像为例，白天识别模型转向夜间图像，若只考察来源知识能否改善夜间识别，重点就是
迁移；若系统随时间接触不同环境，并要求继续保持已学会的白天识别能力，就还要考察持续学习。
若同一系统共同学习物体识别与距离估计，任务集合即使保持固定，也可以采用多任务学习。
若利用多个历史场景反复练习从少量示例适应，再在新场景中检查适应效果，则体现了本节所述的
元学习思路。区别落在知识怎样组织和怎样评价，不能只看是否出现新数据或多个任务。

即使都归入跨域迁移，仍须进一步检查目标信息：提前利用夜间数据对应领域自适应；训练与选择
模型时完全见不到夜间数据、随后直接预测，对应领域泛化；部署后利用夜间输入继续更新，
则涉及测试时适应。数据和任务是否变化之外，\emph{何时能访问哪些信息}也是区分设定的必要条件。

这些关系也说明，知识既可以跨任务共享，也可以沿时间积累。一组任务共同学到的表示可以成为
之后迁移的起点，元学习形成的适应方式也可以用于持续到来的任务。能否组合，不仅取决于任务
是否相关，还取决于当前能否访问所需数据，以及评价是否覆盖真正要保留和获得的能力。

\section{其他范式}
\label{sec:other-learning-settings}

即使反馈来源和任务目标已经确定，学习仍可能受到发生方式与环境的约束。数据可能逐条到来，
要求模型不断预测和更新；也可能分散在不同设备上，只能通过协作利用。这些条件分别涉及
学习的时序，以及数据的存放与协作方式。本节将在线学习与联邦学习归为其他范式，它们可以
与前面的学习设定组合。

\subsection{在线学习}
\label{sec:online-learning}

如果预测请求和反馈持续到来，等到所有数据收齐再训练就不合适了。\term{在线学习}
（online learning）让学习者在经验依次到来的过程中作出预测或决策，并利用随后获得的反馈
更新。它的典型过程是用当前模型预测，取得反馈后调整，再处理后续样本；与之相对，
\term{批量学习}（batch learning）先利用一批给定数据训练，再使用得到的模型。

广告点击率预测提供了一个实际例子：系统根据广告内容和展示场景预测点击，之后观察用户
是否点击，用不断到来的反馈更新模型\scite{mcmahan2013}。
这里点击结果提供监督目标，“在线”则说明这些反馈按什么时序参与学习。反馈可能延迟到达，
更新也可以按小批次进行；在线学习既不要求必须联网，也不意味着部署一个提供实时预测的模型
就已经在在线学习。若模型始终不利用新反馈更新，它只是持续提供预测。

在线学习与持续学习会重叠，但关注点不同。即使新样本始终来自同一分布，也可以依次预测和
更新；在线学习本身不以任务变化或保留一组旧任务为必要条件。持续学习强调长期适应与能力
保留，也可以按阶段接收一批数据再更新。一个逐条学习新反馈、同时检查旧能力的系统，则同时
具有这两种特征。

数据按时间到来时，还可能发生\term{概念漂移}（concept drift），即产生数据的统计规律随
时间改变。例如，同样的广告和用户特征，过去容易带来点击，现在却不再有效；也可能是常见
用户或广告类型的比例改变。前一种变化影响给定输入后的目标规律，后一种改变输入分布。
这类变化可能使旧经验不再充分代表当前情况，需要检测变化或调整不同时间经验的权重
\scite{bifet2007}。在线更新提供了利用新经验的时序安排，但不会自动消除漂移的影响；它还要
应对反馈延迟与新样本不足。第~\ref{sec:distribution-shift}节将进一步区分分布中哪些部分发生
变化，而这里强调这些变化沿时间出现。

评价在线学习时，应按实际可用信息的时间顺序记录预测：先保存模型尚未见到答案时的结果，
再在反馈到达后更新。若用已经参与更新的样本重新预测，就无法反映系统当时面对新数据的表现。
还需要说明反馈延迟、允许保留的数据量和更新预算，因为这些条件决定了模型何时能够利用经验。

\subsection{联邦学习}
\label{sec:federated-learning}

手机键盘要预测用户接下来会输入什么，需要学习多样的用词和表达方式。每台手机只有一部分
输入记录，其他手机虽有相关数据，却不便把原始文本都集中到一处。这里缺少的未必是全局
数据量，而是一种共同利用分散数据的办法。

\term{联邦学习}（federated learning）让不同设备或机构保留本地数据，通过交换模型更新
等信息协作学习。手机键盘的相关研究就在设备端训练下一词预测模型，再汇总各设备的学习
结果以改善预测\scite{hard2018}。下一词目标来自文本本身，本地可以采用自监督学习；
联邦学习则回答这些本地经验怎样共同参与训练。两者描述同一过程的不同方面。

一种常见的中心协调过程是：服务器向一批设备发送当前模型，设备用本地数据继续学习，再把
学到的模型变化发送给服务器。这里交换的“模型更新”可以是模型参数或参数的改变量，服务器
将各方结果汇总，形成下一轮使用的共享模型。\term{联邦平均}（federated averaging, FedAvg）
就是这种组织方式的代表方法，它通过加权平均各设备训练后的模型参数来形成共享模型
\scite{mcmahan2017}。
图~\ref{fig:federated-learning-process}展示了模型下发、本地学习和更新汇总的循环。

\begin{figure*}[t]
  \centering
  % 三个本地域的浅灰底板只表示数据边界；数据不出域，参数通信出域。
% 外层蓝色表示共享模型下发，红色表示本地参数上报；网络内蓝红黄只区分层。
\begingroup%
% 第1—3章局部矢量图元：单色黄、双色蓝红、三色蓝红黄。
\providecommand{\BasicsFiveSetup}{%
 \colorlet{B5Blue}{FigureBlue}\colorlet{B5Coral}{FigureRed}%
 \colorlet{B5Yellow}{FigureYellow}\definecolor{B5Gray}{HTML}{4C4D4F}%
 \colorlet{B5Ice}{FigureBlueFill}\colorlet{B5Rose}{FigureRedFill}%
 \colorlet{B5Cream}{FigureYellowFill}%
 \tikzset{b5 picture/.style={x=1mm,y=1mm,font=\figurefont\mdseries\fontsize{8.5}{10.5}\selectfont,text=B5Gray,line width=1pt},%
 b5 node/.style={draw=B5Gray,fill=white,line width=1pt,rounded corners=1.5mm,inner sep=3pt,font=\figurefont\mdseries\fontsize{8.5}{10.5}\selectfont,text=B5Gray,align=center},%
 b5 flow/.style={draw=B5Gray,line width=1.1pt,-{Stealth[length=2.2mm,width=1.4mm]}},%
 b5 link/.style={draw=B5Gray,line width=.7pt},%
 b5 feedback/.style={b5 flow,draw=B5Coral,dash pattern=on 2.5mm off 1.5mm}}}%
% 3—3—1网络仅为图标，不指定真实架构；各层用实心蓝、红、黄，连线中性灰。
\providecommand{\BasicsFiveNet}[3]{\begin{scope}[shift={(#1,#2)},scale=#3]%
 \foreach \a in {-3,0,3}{\foreach \b in {-3,0,3}{\draw[b5 link] (0,\a)--(5,\b);}}%
 \foreach \a in {-3,0,3}{\draw[b5 link] (5,\a)--(10,0);}%
 \foreach \a in {-3,0,3}{%
  \draw[draw=B5Gray,line width=.8pt,fill=B5Blue] (0,\a) circle (1.1);%
  \draw[draw=B5Gray,line width=.8pt,fill=B5Coral] (5,\a) circle (1.1);}%
 \draw[draw=B5Gray,line width=.8pt,fill=B5Yellow] (10,0) circle (1.1);%
\end{scope}}%
% 鸟与猫保留教学对象，不生成或模拟真实数据样本。
\providecommand{\BasicsFiveBird}[4]{\begin{scope}[shift={(#1,#2)},scale=#3]%
 \draw[draw=B5Gray,fill=#4,line width=.8pt,line join=round]%
  (-6,-1)--(-10,-2)--(-6,2).. controls (-3,3) and (-2,6) .. (1,6)%
  .. controls (4,6) and (5,4) .. (4,2).. controls (4,-1) and (-1,-3) .. (-6,-1);%
 \draw[draw=B5Gray,fill=white,line width=.7pt] (4,3)--(6,3.8)--(4,4.1)--cycle;%
 \fill[B5Gray] (2.7,4.3) circle (.35);%
 \draw[b5 link] (-5,1).. controls (-1,4) and (1,1) .. (-3,-.5);%
 \draw[b5 link] (-2,-2)--(-2,-4)--(-.7,-4);\draw[b5 link] (0,-1.5)--(0,-4)--(1.3,-4);%
\end{scope}}%
\providecommand{\BasicsFiveCat}[3]{\begin{scope}[shift={(#1,#2)},scale=#3]%
 \draw[draw=B5Gray,fill=white,line width=.8pt] (0,-1) ellipse (3.3 and 4.3);%
 \draw[draw=B5Gray,fill=white,line width=.8pt,line join=round]%
  (-3,3)--(-3,7)--(-1.2,5.8).. controls (0,6.5) and (1,6.5) .. (2,5.8)%
  --(3.3,7)--(3.3,3).. controls (2,1.3) and (-1.5,1.3) .. (-3,3);%
 \fill[B5Gray] (-1.3,4) circle (.3);\fill[B5Gray] (1.5,4) circle (.3);%
 \draw[b5 link] (0,3.4)--(0,2.7);\draw[b5 link] (-2,-2)--(-2,-5);\draw[b5 link] (2,-2)--(2,-5);%
 \draw[draw=B5Gray,line width=.8pt] (-2.5,-3).. controls (-8,-5) and (-8,-.5) .. (-5,-1);%
\end{scope}}%
\providecommand{\BasicsFiveRobot}[3]{\begin{scope}[shift={(#1,#2)},scale=#3]%
 \fill[B5Gray,rounded corners=.4mm] (-2.7,-1.6) rectangle (-1.7,1.6);%
 \fill[B5Gray,rounded corners=.4mm] (1.7,-1.6) rectangle (2.7,1.6);%
 \draw[draw=B5Gray,fill=white,line width=1pt,rounded corners=.7mm] (-1.8,-2.3) rectangle (1.8,2.3);%
 \draw[draw=B5Gray,fill=B5Yellow,line width=.8pt] (0,.6) circle (.65);%
\end{scope}}%
\BasicsFiveSetup%
\begin{tikzpicture}[b5 picture]%
 \path[use as bounding box] (0,-3) rectangle (168,100);%
 \foreach \left/\ii in {0/1,55/2,110/3}{%
  \draw[draw=B5Gray,fill=FigurePanel,line width=1pt,rounded corners=2mm] (\left,20.5) rectangle (\left+50,65);%
  \node[anchor=west] at (\left+2,68) {参与方\ii};%
  \ifnum\ii=1%
   \draw[draw=B5Gray,fill=white,line width=.8pt] (\left+5,28) rectangle (\left+13,34);%
   \draw[draw=B5Gray,fill=FigurePanel,line width=.7pt] (\left+6,29) rectangle (\left+12,33);%
   \draw[draw=B5Gray,fill=white,line width=.8pt] (\left+5,28)--(\left+3,26.5)--(\left+15,26.5)--(\left+13,28);%
  \else\ifnum\ii=2%
   \draw[draw=B5Gray,fill=white,line width=.8pt] (\left+4,29) rectangle (\left+14,35);%
   \draw[draw=B5Gray,fill=FigurePanel,line width=.7pt] (\left+5,30) rectangle (\left+13,34);%
   \draw[b5 link] (\left+9,29)--(\left+9,27)--(\left+6,27)--(\left+12,27);%
  \else%
   \draw[draw=B5Gray,fill=white,line width=.8pt,rounded corners=.7mm] (\left+7,26.5) rectangle (\left+12,35);%
   \draw[draw=B5Gray,fill=FigurePanel,line width=.7pt] (\left+7.7,28.5) rectangle (\left+11.3,33.5);%
   \draw[draw=B5Gray,line width=.6pt] (\left+9.5,27.5) circle (.3);%
  \fi\fi%
  \draw[draw=B5Gray,fill=white,line width=1pt] (\left+4,48)--(\left+4,56) arc[start angle=180,end angle=360,x radius=5,y radius=1.6]--(\left+14,48) arc[start angle=0,end angle=-180,x radius=5,y radius=1.6];%
  \draw[draw=B5Gray,fill=white,line width=1pt] (\left+9,56) ellipse (5 and 1.6);%
  \node at (\left+9,51) {$D_{\ii}$};%
  \node[b5 node,minimum width=26mm,minimum height=16mm] (local\ii) at (\left+35,52) {};%
  \node at (\left+35,57) {本地学习};\BasicsFiveNet{\left+29}{50}{1.05}%
  \draw[b5 flow] (\left+14,52)--(local\ii.west);%
  \node[b5 node,minimum width=26mm,minimum height=20mm] (update\ii) at (\left+35,32) {};%
  \node at (\left+35,37) {$\bm\theta_{\ii}^{(t+1)}$};\BasicsFiveNet{\left+29}{29.5}{1.05}%
  \draw[b5 flow,draw=B5Coral] (local\ii.south)--(update\ii.north);}%
 \node[b5 node,fill=B5Ice,minimum width=38mm,minimum height=22mm] (shared) at (84,87) {};%
 \node at (84,94) {共享模型 $\bm\theta^{(t)}$};\BasicsFiveNet{78}{83}{1.2}%
 \draw[b5 flow,draw=B5Blue,rounded corners=1.5mm] ([xshift=-13mm]shared.south)--(71,73)--(35,73)--(local1.north);%
 \draw[b5 flow,draw=B5Blue,rounded corners=1.5mm] (shared.south)--(84,73)--(90,73)--(local2.north);%
 \draw[b5 flow,draw=B5Blue,rounded corners=1.5mm] ([xshift=13mm]shared.south)--(97,73)--(145,73)--(local3.north);%
 \node[text=B5Blue] at (39,77) {模型下发};%
 \node[text=B5Blue,anchor=west] at (92,70) {模型下发};%
 \node[text=B5Blue] at (136,77) {模型下发};%
 \node[b5 node,fill=B5Rose,minimum width=36mm,minimum height=20mm] (aggregate) at (84,8) {};%
 \node at (84,14.5) {加权聚合 $\bm\theta^{(t+1)}$};\BasicsFiveNet{78.75}{4.5}{1.05}%
 \draw[b5 flow,draw=B5Coral,rounded corners=1.5mm] (update1.south)--(35,8)--(aggregate.west);%
 \draw[b5 flow,draw=B5Coral,rounded corners=1mm] (update2.south)--(90,21)--(84,21)--(aggregate.north);%
 \draw[b5 flow,draw=B5Coral,rounded corners=1.5mm] (update3.south)--(145,8)--(aggregate.east);%
 \node[text=B5Coral] at (48,14) {模型更新};%
 \node[text=B5Coral,anchor=east] at (65,17) {模型更新};%
 \node[text=B5Coral] at (129,14) {模型更新};%
 \draw[b5 flow,draw=B5Blue,rounded corners=2mm] ([yshift=-6mm]aggregate.east)--(165,2)--(165,87)--(shared.east);%
 \node[text=B5Gray] at (136,91) {下一轮};%
\end{tikzpicture}%
\endgroup%

  \caption{中心协调式联邦学习的一轮通信。协调方向三个示意参与方下发同一组共享参数$\bm\theta^{(t)}$；各方使用边界内的本地记录$D_i$学习，上传本地模型更新。加权聚合形成下一版本$\bm\theta^{(t+1)}$，再统一下发。记录图元与参与方数量为教学示意；原始记录留在本地，但模型更新仍可能泄露信息，联邦组织方式本身不提供隐私保证。}
  \label{fig:federated-learning-process}
\end{figure*}
\FloatBarrier

这个循环要面对实际设备的差异：有的设备数据多，有的设备算力弱，还可能有设备中途离线。
模型反复下发和上传也需要通信成本。原始数据虽然留在本地，模型更新仍可能透露本地信息，
所以隐私保护还需要额外措施，例如限制谁能访问单个设备的更新。协作训练中的模型聚合、
通信与部署问题由
\BookPartRef{part:systems}{第六卷“系统”的“机器学习系统”部分}展开，隐私保护方法则在
第\ref{part:trustworthy-ml}部分“机器学习可信性”中讨论。

联邦学习可以缓解单个参与方的数据不足，但不同用户的用词习惯仍可能不同。它可以与迁移学习
结合，让某台设备在共享模型的基础上适应本地数据；也可以采用多任务学习，让不同参与方的
任务共享信息，同时保留各自的差异。若用户习惯又随时间变化，还会出现持续学习问题。因此，
协作利用分散数据、适应各方差异和保留旧能力，可能是同一个系统需要同时解决的问题。

\section{范式的组合}
\label{sec:combining-learning-paradigms}

前面的范式分别约束反馈来源、标签利用、知识复用和学习环境。同一个系统可能同时面对几类
问题，因此可以在不同阶段采用不同范式，也可以在同一阶段利用多种信号。理解组合时，要
分别追踪各阶段的数据、目标与反馈，并说明哪些知识被传给下一阶段。

预训练与下游学习是一种常见组合。模型先在大量无标签文本上用自监督目标学习表示，再用带
任务标签的数据进行监督学习。前一阶段提供可复用的知识，后一阶段把它调整到问答、分类或
其他具体任务。半监督学习则在同一目标任务中共同使用有标签和无标签数据，有些方法还会附加
自监督目标。两者都包含多种信号，但组织方式和问题设定并不相同。

2016年的AlphaGo提供了棋谱与胜负反馈结合的实例。它先让负责选择落子的策略网络从专家
棋谱中的局面与落子学习，得到初始策略，再通过自我对弈，根据最终胜负继续改进策略
\scite{silver2016}。从棋谱到自我对弈，传递的是已经学到的模型；后续反馈则从专家动作
变为对局结果。图~\ref{fig:paradigm-combination-alphago}只呈现这条学习主线，完整系统还
包含搜索等组件。这里的自我对弈属于强化学习过程，不能因名称中有“自我”就归为自监督学习。

\begin{figure}[!htb]
  \centering
  \bookdraftgraphic{figures/scenes/v1r3-combination-alphago.png}
  \caption{AlphaGo学习主线的简化示意。人类棋谱为监督学习阶段提供局面与专家落子；得到初始策略后，
  自我对弈产生新棋局与胜负反馈，用于强化学习阶段。完整系统还包含搜索等组件。}
  \label{fig:paradigm-combination-alphago}
\end{figure}
\FloatBarrier

还有些行为难以直接打分，却比较容易比较好坏。例如，我们未必能为机器人的每一步动作写出
奖励，但可以比较两段行为哪一段更符合要求。\term{偏好学习}（preference learning）利用
的就是这类相对判断。它与强化学习的一种组合方式是向人展示两段行为，请人选出更满意的一段，
用比较结果训练一个评价模型，再用该模型提供的奖励训练智能体。相关研究曾将这种方法用于
Atari游戏和模拟机器人任务
\scite{christiano2017}。这里，人提供的是“哪段行为更好”的比较，而不是每一步该怎么做的
动作示范；评价模型把这种偏好转成奖励，强化学习再利用奖励改进行动。

组合也可以发生在同一个训练过程中。以一个假设的花卉识别系统为例：不同机构各自保留照片，
通过联邦学习协作训练；各机构从本地无标签照片构造不同视图，利用自监督对比目标学习表示。
其中，联邦学习规定数据怎样分布与协作，自监督学习说明匹配信号从哪里来，对比学习说明怎样
利用这些关系形成目标。三种描述各有作用，并不相互替代。

若还能够请专家标注部分照片，系统可以主动选择需要询问的样本，再结合其余无标签照片进行
半监督学习。已有通用图像表示可以迁移到花卉识别；随着新物种陆续加入，系统又需要检查持续
学习中的旧类别保留。这是为说明组合而构造的教学情景，实际方案只需采用问题所要求的部分。
标签选择不能代替目标任务的学习，协作训练也不会自动解决来源差异或遗忘。

\section{本章小结}
\label{sec:machine-learning-paradigms-summary}

学习不必依赖完整标签；关键是先确定系统能从数据、交互或已有任务中获得什么反馈，再用相应
范式组织训练和评价。任务、模型与学习范式分别说明要完成什么、怎样表示规律，以及依据什么
经验学习这些规律。监督学习利用输入及对应目标，无监督学习从数据中发现结构或分布规律，
强化学习则根据行动带来的奖励改进决策。判断反馈的来源，比模型采用什么名称更能说明学习
是怎样发生的。

回到章首的标签不足问题，学习可以利用未标注数据补充已有监督，从数据自身构造目标，接受
较弱但成本较低的监督，或主动获取关键标签。对比学习进一步说明怎样用样本关系形成表示
学习目标。这些办法能否发挥作用，取决于额外信号是否可靠、是否与目标任务有关；减少人工
标签并不保证模型自动获得所需能力。

数据或任务变化后，还要明确知识服务的对象与需要保留的能力。持续学习考察顺序更新中的
适应与保留，多任务学习考察相关任务间的共享，迁移学习考察来源知识对目标问题的帮助，
元学习则用历史任务改进后续学习过程。在线学习与联邦学习进一步限定更新时序和协作环境。
目标域信息何时可用，又区分了领域自适应、领域泛化与测试时适应。
这些维度可以组合，因此分析一个系统时，应分别说明数据、反馈、任务关系和信息可用条件，
并按最终需要获得或保留的能力评价结果。

\paragraph{延伸阅读}

若要深入理解交互决策，可从Sutton与Barto的强化学习教材继续；若关注标签质量和任务适应，
可分别阅读弱监督学习综述与元学习综述；联邦平均的原始研究则提供了分散数据协作训练的
具体入口\scite{sutton2018,zhou2018,hospedales2022,mcmahan2017}。结合
第~\ref{chap:machine-learning-models}章，我们已经能分别说明怎样表示规律，以及怎样利用
经验学习这些规律。第\ref{part:learning-theory}部分将进一步讨论有限经验在什么条件下能够
支持泛化。\BookVolumeRef{vol:paradigms}{第四卷}“范式”将围绕强化学习、知识的高效利用、知识的演进
与迁移展开具体方法，协作训练与部署约束则衔接系统卷。
